\documentclass[aps,preprint]{revtex4}%
\usepackage{amssymb}
\usepackage{amsfonts}
\usepackage{amsmath}
\usepackage{graphicx}%
\setcounter{MaxMatrixCols}{30}
%TCIDATA{OutputFilter=latex2.dll}
%TCIDATA{Version=5.50.0.2960}
%TCIDATA{CSTFile=revtex4.cst}
%TCIDATA{Created=Friday, February 06, 2015 09:42:56}
%TCIDATA{LastRevised=Sunday, May 24, 2015 08:50:41}
%TCIDATA{<META NAME="GraphicsSave" CONTENT="32">}
%TCIDATA{<META NAME="SaveForMode" CONTENT="1">}
%TCIDATA{BibliographyScheme=Manual}
%TCIDATA{<META NAME="DocumentShell" CONTENT="Articles\SW\REVTeX 4">}
%BeginMSIPreambleData
\providecommand{\U}[1]{\protect\rule{.1in}{.1in}}
%EndMSIPreambleData
\newtheorem{theorem}{Theorem}
\newtheorem{acknowledgement}[theorem]{Acknowledgement}
\newtheorem{algorithm}[theorem]{Algorithm}
\newtheorem{axiom}[theorem]{Axiom}
\newtheorem{claim}[theorem]{Claim}
\newtheorem{conclusion}[theorem]{Conclusion}
\newtheorem{condition}[theorem]{Condition}
\newtheorem{conjecture}[theorem]{Conjecture}
\newtheorem{corollary}[theorem]{Corollary}
\newtheorem{criterion}[theorem]{Criterion}
\newtheorem{definition}[theorem]{Definition}
\newtheorem{example}[theorem]{Example}
\newtheorem{exercise}[theorem]{Exercise}
\newtheorem{lemma}[theorem]{Lemma}
\newtheorem{notation}[theorem]{Notation}
\newtheorem{problem}[theorem]{Problem}
\newtheorem{proposition}[theorem]{Proposition}
\newtheorem{remark}[theorem]{Remark}
\newtheorem{solution}[theorem]{Solution}
\newtheorem{summary}[theorem]{Summary}
\newenvironment{proof}[1][Proof]{\noindent\textbf{#1.} }{\ \rule{0.5em}{0.5em}}
\begin{document}
\preprint{HEP/123-qed}
\title[Short title for running header]{REV\TeX4}
\author{First Author}
\affiliation{My Institution}
\author{Second Author}
\affiliation{My Institution}
\author{Third Author}
\affiliation{Other Institution}
\keywords{one two three}
\pacs{PACS number}

\begin{abstract}
Shell document for REV\TeX{} 4.

\end{abstract}
\volumeyear{year}
\volumenumber{number}
\issuenumber{number}
\eid{identifier}
\date[Date text]{date}
\received[Received text]{date}

\revised[Revised text]{date}

\accepted[Accepted text]{date}

\published[Published text]{date}

\startpage{101}
\endpage{102}
\maketitle
\tableofcontents


\section{Canonical Form of a Real Antisymmetric Matrix\label{realanti}}

A real antisymmetric bilinear form
\begin{equation}
g:\mathbb{R}^{r}\times\mathbb{R}^{r}\rightarrow\mathbb{R}%
\end{equation}
can be expressed as
\begin{equation}
g\left(  a,b\right)  =\left\langle g\left\vert a\wedge b\right.  \right\rangle
=\left\langle a\left\vert Ab\right.  \right\rangle =\mathbf{a}^{t}%
\mathbb{A}\mathbf{b},
\end{equation}
where
\begin{equation}
\left\vert g\right\rangle =\frac{1}{2}\sum_{1\leq j_{1},j_{2}\leq r}%
A_{j_{1}j_{2}}\left\vert e_{j_{1}}e_{j_{2}}\right\rangle ;\quad A_{j_{1}j_{2}%
}=-A_{j_{2}j_{1}}%
\end{equation}
and
\begin{equation}
\left\vert e_{j_{1}}e_{j_{2}}\right\rangle =\frac{1}{\sqrt{2}}\left(
\left\vert e_{j_{1}}\otimes e_{j_{2}}\right\rangle -\left\vert e_{j_{2}%
}\otimes e_{j_{1}}\right\rangle \right)  .
\end{equation}


The preceding expressions can be formulated in a canonical way by observing
that a real antisymmetric\emph{ }matrix $\mathbb{A}=\mathbb{A}^{\ast
}=-\mathbb{A}^{t}$ can be converted into a \emph{purely imaginary hermitian}
\emph{ }antisymmetric matrix, $\mathbb{H}$, $\lceil\mathcal{B}\left(
\mathbb{R}^{2s}\right)  \hookrightarrow\mathcal{B}\left(  \mathbb{C}%
^{2s}\right)  \rfloor$ with the properties
\begin{align}
\mathbb{H}  &  =i\mathbb{A}\nonumber\\
\mathbb{H}^{\dagger}  &  =-i\mathbb{A}^{t}=i\mathbb{A}=\mathbb{H}\nonumber\\
\mathbb{H}^{\ast}  &  =-i\mathbb{A}=i\mathbb{A}^{t}\,=\mathbb{H}%
^{t}=-\mathbb{H}%
\end{align}
and
\begin{equation}
\mathbb{H}=i\mathbb{A}=\sum_{1\leq j\leq s}\alpha_{j}\left\{  \left\vert
\mathbf{v}_{j}\right\rangle \left\langle \mathbf{v}_{j}\right\vert -\left\vert
\mathbf{v}_{j}^{\ast}\right\rangle \left\langle \mathbf{v}_{j}^{\ast
}\right\vert \right\}  ;\qquad\alpha_{j}\in\mathbb{R}^{+}, \label{sa}%
\end{equation}
where $\left\{  \left\vert \mathbf{v}_{j}\right\rangle ,\left\vert
\mathbf{v}_{j}^{\ast}\right\rangle ;1\leq j\leq r\right\}  $ is an orthonormal basis.

$\lceil$%
\begin{equation}
v_{j}\left(  e_{1},\cdots,e_{r}\right)  =\left(  e_{1},\cdots,e_{r}\right)
\left\vert \mathbf{v}_{j}\right\rangle =\left(  e_{1},\cdots,e_{r}\right)
R_{e}\left(  v_{j}\right)
\end{equation}
$\rfloor$

Eq. $\left(  \ref{sa}\right)  $ is shown by
\begin{equation}
\mathbb{H}\left\vert \mathbf{v}\right\rangle =\alpha\left\vert \mathbf{v}%
\right\rangle \Rightarrow\mathbb{H}^{\ast}\left\vert \mathbf{v}^{\ast
}\right\rangle =\alpha\left\vert \mathbf{v}^{\ast}\right\rangle \Rightarrow
\mathbb{H}^{t}\left\vert \mathbf{v}^{\ast}\right\rangle =\alpha\left\vert
\mathbf{v}^{\ast}\right\rangle \Rightarrow\mathbb{H}\left\vert \mathbf{v}%
^{\ast}\right\rangle =-\alpha\left\vert \mathbf{v}^{\ast}\right\rangle ,
\end{equation}
If $\alpha\neq0$ then $\left\vert \mathbf{v}^{\ast}\right\rangle
\bot\left\vert \mathbf{v}\right\rangle ,$ if $\alpha=0$ then $\left\vert
\mathbf{v}^{\ast}\right\rangle $ is degenerate with $\left\vert \mathbf{v}%
\right\rangle $ and one can always choose an imaginary basis for the kernel of
$\mathbb{H}$, which could have odd dimension. Thus in the case $r\neq2s$ there
is always at least one zero eigenvalue of $\mathbb{H}$ i.e. the dimension of
Ker$\left\{  \mathbb{H}\right\}  \geq1.$

$\lceil$%
\begin{align}
\mathbb{H}\left\vert \mathbf{v}\right\rangle  &  =\alpha\left\vert
\mathbf{v}\right\rangle \\
\mathbb{H}\left\vert \mathbf{v}^{\ast}\right\rangle  &  =-\alpha\left\vert
\mathbf{v}^{\ast}\right\rangle .
\end{align}
$\rfloor$

Eq. $\left(  \ref{sa}\right)  $ can also be expressed as
\begin{equation}
\mathbb{H}=\mathbb{V}\left(
\begin{array}
[c]{cc}%
\boldsymbol{\alpha} & \mathbb{O}\\
\mathbb{O} & -\boldsymbol{\alpha}%
\end{array}
\right)  \mathbb{V}^{\dagger}%
\end{equation}
leading to
\begin{equation}
\mathbb{A}=-i\mathbb{H=}-i\mathbb{V}\left(
\begin{array}
[c]{cc}%
\boldsymbol{\alpha} & \mathbb{O}\\
\mathbb{O} & -\boldsymbol{\alpha}%
\end{array}
\right)  \mathbb{V}^{\dagger}.
\end{equation}


$\lceil$Let
\begin{equation}
\mathbb{A}=\left(
\begin{array}
[c]{cc}%
\mathbb{O} & -\mathbb{I}_{2}\\
\mathbb{I}_{2} & \mathbb{O}%
\end{array}
\right)  =\mathbb{J}%
\end{equation}
then
\begin{equation}
\mathbb{H}=i\mathbb{A}=\left\vert \mathbf{v}_{1}\right\rangle \left\langle
\mathbf{v}_{1}\right\vert +\left\vert \mathbf{v}_{2}\right\rangle \left\langle
\mathbf{v}_{2}\right\vert -\left(  \left\vert \mathbf{v}_{3}^{\ast
}\right\rangle \left\langle \mathbf{v}_{3}^{\ast}\right\vert +\left\vert
\mathbf{v}_{4}^{\ast}\right\rangle \left\langle \mathbf{v}_{4}^{\ast
}\right\vert \right)  =i\left(
\begin{array}
[c]{cc}%
\mathbb{O} & -\mathbb{I}_{2}\\
\mathbb{I}_{2} & \mathbb{O}%
\end{array}
\right)  =i\mathbb{J}%
\end{equation}
and
\begin{equation}
\mathbb{J=-}ii\mathbb{H}.
\end{equation}
Note that complexification does not have to be antisymmetric$\rfloor$

\subsubsection{$\lceil$ Notation}

Vectors in bold face are numerical representations wrt a basis $\left\{
e_{1},\cdots,e_{2s}\right\}  $ so the row super vector of numerical vectors
$\left\vert \mathbf{v}_{1},\cdots,\mathbf{v}_{2s}\right\rangle $ can be viewed
as the matrix
\begin{equation}
\mathbb{V}=\left(  \left.  \mathbf{v}_{1}\right\vert ,\cdots,\left\vert
\mathbf{v}_{2s}\right.  \right)
\end{equation}
whose columns are made up of the numerical vectors $\left\{  \mathbf{v}%
_{1},\cdots,\mathbf{v}_{2s}\right\}  $ and
\begin{align}
&  V\left(  e_{1},\cdots,e_{2s}\right)  =\left(  e_{1},\cdots,e_{2s}\right)
\mathbb{V}\\
\mathbb{V}  &  =\left(  \mathbf{e}_{1},\cdots,\mathbf{e}_{2s}\right)
\mathbb{V}=\left(  \left.  \mathbf{e}_{1}\right\vert ,\cdots,\left\vert
\mathbf{e}_{2s}\right.  \right)  \mathbb{V}=\left(
\begin{array}
[c]{ccc}%
1 & 0 & 0\\
0 & \ddots & 0\\
0 & 0 & 1
\end{array}
\right)  \mathbb{V}.
\end{align}
$\rfloor$

It is important to note that the orthonormal basis eigenvectors of
$\mathbb{H}$ are defined only up to overall phase factors $e^{i\varphi_{j}}$
which \emph{do not} preserve the complex conjugate relationship between
eigenvectors associated with paired eigenvalues of opposite sign, although one
can always find complex conjugate eigenvector pairs. These different bases are
related to each other by
\begin{equation}
\left\vert \mathbf{v}_{1}^{\prime},\cdots,\mathbf{v}_{2s}^{\prime
}\right\rangle =\left\vert \mathbf{v}_{1},\cdots,\mathbf{v}_{s},\mathbf{v}%
_{1}^{\ast},\cdots,\mathbf{v}_{s}^{\ast}\right\rangle \left(
\begin{array}
[c]{ccc}%
e^{i\varphi_{1}} & \mathbf{0} & \mathbf{0}\\
\mathbf{0} & \ddots & \mathbf{0}\\
\mathbf{0} & \mathbf{0} & e^{i\varphi_{2s}}%
\end{array}
\right)  .
\end{equation}


$\lceil$%
\begin{equation}
\left\vert \mathbf{v}_{1}^{\prime},\cdots,\mathbf{v}_{4}^{\prime}\right\rangle
=\left\vert \mathbf{v}_{1},\mathbf{v}_{2},\mathbf{v}_{1}^{\ast},\mathbf{v}%
_{2}^{\ast}\right\rangle \left(
\begin{array}
[c]{ccc}%
e^{i\varphi_{1}} & \mathbf{0} & \mathbf{0}\\
\mathbf{0} & \ddots & \mathbf{0}\\
\mathbf{0} & \mathbf{0} & e^{i\varphi_{4}}%
\end{array}
\right)  .
\end{equation}


$\rfloor$

Using a basis with complex conjugate symmetry we obtain
\begin{equation}
\mathbb{A}=-\sum_{1\leq j\leq s}i\alpha_{j}\left\{  \left\vert \mathbf{v}%
_{j}\right\rangle \left\langle \mathbf{v}_{j}\right\vert -\left\vert
\mathbf{v}_{j}^{\ast}\right\rangle \left\langle \mathbf{v}_{j}^{\ast
}\right\vert \right\}  ;\qquad\alpha_{i}\in\mathbb{R}.
\end{equation}
$\lceil$%
\begin{equation}
\mathbb{A}=-i\left\{  \left\vert \mathbf{v}_{1}\right\rangle \left\langle
\mathbf{v}_{1}\right\vert +\left\vert \mathbf{v}_{2}\right\rangle \left\langle
\mathbf{v}_{2}\right\vert -\left(  \left\vert \mathbf{v}_{3}^{\ast
}\right\rangle \left\langle \mathbf{v}_{3}^{\ast}\right\vert +\left\vert
\mathbf{v}_{4}^{\ast}\right\rangle \left\langle \mathbf{v}_{4}^{\ast
}\right\vert \right)  \right\}  .
\end{equation}


$\rfloor$

We can then define a \emph{real} basis $\left\{  \left\vert \mathbf{s}%
_{j}\right\rangle ;1\leq j\leq r\right\}  $ i.e. $\left\vert \mathbf{s}%
_{j}\right\rangle =\left\vert \mathbf{s}_{j}^{\ast}\right\rangle $ as
\begin{align}
\left\vert \mathbf{s}_{j}\right\rangle  &  =\frac{1}{\sqrt{2}}\left(
\left\vert \mathbf{v}_{j}\right\rangle +\left\vert \mathbf{v}_{j}^{\ast
}\right\rangle \right)  =\sqrt{2}\operatorname{Re}\left\vert \mathbf{v}%
_{j}\right\rangle \nonumber\\
\left\vert \mathbf{s}_{j+\mathbf{s}}\right\rangle  &  =\frac{-i}{\sqrt{2}%
}\left(  \left(  \left\vert \mathbf{v}_{j}\right\rangle -\left\vert
\mathbf{v}_{j}^{\ast}\right\rangle \right)  \right)  =\sqrt{2}%
\operatorname{Im}\left\vert \mathbf{v}_{j}\right\rangle ,
\end{align}
i.e.
\begin{equation}
\left\vert \mathbf{s}_{1},\cdots,\mathbf{s}_{\mathbf{s}},\mathbf{s}%
_{\mathbf{s}+1},\cdots,\mathbf{s}_{\mathbf{s}+\mathbf{s}}\right\rangle
=\left\vert \mathbf{v}_{1},\cdots,\mathbf{v}_{\mathbf{s}},\mathbf{v}_{1}%
^{\ast},\cdots,\mathbf{v}_{\mathbf{s}}^{\ast}\right\rangle \frac{1}{\sqrt{2}%
}\left(
\begin{array}
[c]{cc}%
\mathbb{I}_{s} & -i\mathbb{I}_{s}\\
\mathbb{I}_{s} & i\mathbb{I}_{s}%
\end{array}
\right)  ,
\end{equation}
which can also be written as
\begin{equation}
\mathbb{O}=\mathbb{V}\frac{1}{\sqrt{2}}\left(
\begin{array}
[c]{cc}%
\mathbb{I}_{s} & -i\mathbb{I}_{s}\\
\mathbb{I}_{s} & i\mathbb{I}_{s}%
\end{array}
\right)  ,
\end{equation}
where
\begin{align}
\mathbb{O}  &  =\left\vert \mathbf{s}_{1},\cdots,\mathbf{s}_{\mathbf{s}%
},\mathbf{s}_{\mathbf{s}+1},\cdots,\mathbf{s}_{\mathbf{s}+\mathbf{s}%
}\right\rangle \\
\mathbb{V}  &  =\left\vert \mathbf{v}_{1},\cdots,\mathbf{v}_{\mathbf{s}%
},\mathbf{v}_{1}^{\ast},\cdots,\mathbf{v}_{\mathbf{s}}^{\ast}\right\rangle
=\left(  \left.  \mathbb{V}_{s}\right\vert \mathbb{V}_{s}^{\ast}\right)  .
\end{align}


$\lceil$%
\begin{align}
\frac{1}{\sqrt{2}}\left(
\begin{array}
[c]{cc}%
\mathbb{I}_{s} & -i\mathbb{I}_{s}\\
\mathbb{I}_{s} & i\mathbb{I}_{s}%
\end{array}
\right)  \frac{1}{\sqrt{2}}\left(
\begin{array}
[c]{cc}%
\mathbb{I}_{s} & -i\mathbb{I}_{s}\\
\mathbb{I}_{s} & i\mathbb{I}_{s}%
\end{array}
\right)  ^{\dagger}  &  =\frac{1}{2}\left(
\begin{array}
[c]{cc}%
\mathbb{I}_{s} & -i\mathbb{I}_{s}\\
\mathbb{I}_{s} & i\mathbb{I}_{s}%
\end{array}
\right)  \left(
\begin{array}
[c]{cc}%
\mathbb{I}_{s} & \mathbb{I}_{s}\\
i\mathbb{I}_{s} & -i\mathbb{I}_{s}%
\end{array}
\right)  =\left(
\begin{array}
[c]{cc}%
\mathbb{I}_{s} & \mathbb{O}\\
\mathbb{O} & \mathbb{I}_{s}%
\end{array}
\right) \\
\frac{1}{\sqrt{2}}\left(
\begin{array}
[c]{cc}%
\mathbb{I}_{s} & -i\mathbb{I}_{s}\\
\mathbb{I}_{s} & i\mathbb{I}_{s}%
\end{array}
\right)   &  \in U\left(  2s\right)  ;\,\notin SU\left(  2s\right) \\
\det\left\{  \frac{1}{\sqrt{2}}\left(
\begin{array}
[c]{cc}%
\mathbb{I}_{s} & -i\mathbb{I}_{s}\\
\mathbb{I}_{s} & i\mathbb{I}_{s}%
\end{array}
\right)  \right\}   &  =\left(  \frac{1}{\sqrt{2}}\right)  ^{2}\det\left\{
\left(
\begin{array}
[c]{cc}%
\mathbb{I}_{s} & -i\mathbb{I}_{s}\\
\mathbb{I}_{s} & i\mathbb{I}_{s}%
\end{array}
\right)  \right\}  =\frac{2i}{2}=i=e^{\frac{\pi}{2}i}\neq1.
\end{align}%
\begin{align}
e^{\frac{\pi}{2}i}  &  =i\\
e^{\pi i}  &  =-1\\
e^{\frac{3\pi}{2}i}  &  =-i\\
e^{2\pi i}  &  =1
\end{align}
$\rfloor$

$\lceil$%
\begin{align}
\left\vert \mathbf{s}_{1}\right\rangle  &  =\frac{1}{\sqrt{2}}\left(
\left\vert \mathbf{v}_{1}\right\rangle +\left\vert \mathbf{v}_{1}^{\ast
}\right\rangle \right)  =\sqrt{2}\operatorname{Re}\left\vert \mathbf{v}%
_{1}\right\rangle \nonumber\\
\left\vert \mathbf{s}_{2}\right\rangle  &  =\frac{1}{\sqrt{2}}\left(
\left\vert \mathbf{v}_{2}\right\rangle +\left\vert \mathbf{v}_{2}^{\ast
}\right\rangle \right)  =\sqrt{2}\operatorname{Re}\left\vert \mathbf{v}%
_{2}\right\rangle \\
\left\vert \mathbf{s}_{3}\right\rangle  &  =\frac{-i}{\sqrt{2}}\left(  \left(
\left\vert \mathbf{v}_{1}\right\rangle -\left\vert \mathbf{v}_{1}^{\ast
}\right\rangle \right)  \right)  =\sqrt{2}\operatorname{Im}\left\vert
\mathbf{v}_{1}\right\rangle \\
\left\vert \mathbf{s}_{4}\right\rangle  &  =\frac{-i}{\sqrt{2}}\left(  \left(
\left\vert \mathbf{v}_{2}\right\rangle -\left\vert \mathbf{v}_{2}^{\ast
}\right\rangle \right)  \right)  =\sqrt{2}\operatorname{Im}\left\vert
\mathbf{v}_{2}\right\rangle .
\end{align}


$\rfloor$

so that
\begin{align}
\mathbb{A}\left\{  \frac{1}{\sqrt{2}}\left(  \left\vert \mathbf{v}%
_{j}\right\rangle +\left\vert \mathbf{v}_{j}^{\ast}\right\rangle \right)
\right\}   &  =-i\alpha_{j}\frac{1}{\sqrt{2}}\left(  \left\vert \mathbf{v}%
_{j}\right\rangle -\left\vert \mathbf{v}_{j}^{\ast}\right\rangle \right)
\nonumber\\
&  =\alpha_{j}\left\{  \frac{-i}{\sqrt{2}}\left(  \left(  \left\vert
\mathbf{v}_{j}\right\rangle -\left\vert \mathbf{v}_{j}^{\ast}\right\rangle
\right)  \right)  \right\}
\end{align}
and
\begin{align}
\mathbb{A}\left\{  \frac{-i}{\sqrt{2}}\left(  \left(  \left\vert
\mathbf{v}_{j}\right\rangle -\left\vert \mathbf{v}_{j}^{\ast}\right\rangle
\right)  \right)  \right\}   &  =\frac{-\alpha_{j}}{\sqrt{2}}\left\vert
\mathbf{v}_{j}\right\rangle -\frac{\alpha_{j}}{\sqrt{2}}\left\vert
\mathbf{v}_{j}^{\ast}\right\rangle \nonumber\\
&  =-\alpha_{j}\left\{  \frac{1}{\sqrt{2}}\left(  \left\vert \mathbf{v}%
_{j}\right\rangle +\left\vert \mathbf{v}_{j}^{\ast}\right\rangle \right)
\right\}  .
\end{align}
$\lceil$%
\begin{align}
\mathbb{A}\left\{  \frac{1}{\sqrt{2}}\left(  \left\vert \mathbf{v}%
_{1}\right\rangle +\left\vert \mathbf{v}_{1}^{\ast}\right\rangle \right)
\right\}   &  =\frac{-i}{\sqrt{2}}\left(  \left\vert \mathbf{v}_{1}%
\right\rangle -\left\vert \mathbf{v}_{1}^{\ast}\right\rangle \right)
\nonumber\\
\mathbb{A}\left\{  \frac{1}{\sqrt{2}}\left(  \left\vert \mathbf{v}%
_{2}\right\rangle +\left\vert \mathbf{v}_{2}^{\ast}\right\rangle \right)
\right\}   &  =\frac{-i}{\sqrt{2}}\left(  \left\vert \mathbf{v}_{2}%
\right\rangle -\left\vert \mathbf{v}_{2}^{\ast}\right\rangle \right) \\
\mathbb{A}\left\{  \frac{-i}{\sqrt{2}}\left(  \left\vert \mathbf{v}%
_{1}\right\rangle -\left\vert \mathbf{v}_{1}^{\ast}\right\rangle \right)
\right\}   &  =-\left\{  \frac{1}{\sqrt{2}}\left(  \left\vert \mathbf{v}%
_{1}\right\rangle +\left\vert \mathbf{v}_{1}^{\ast}\right\rangle \right)
\right\} \\
\mathbb{A}\left\{  \frac{-i}{\sqrt{2}}\left(  \left\vert \mathbf{v}%
_{2}\right\rangle -\left\vert \mathbf{v}_{2}^{\ast}\right\rangle \right)
\right\}   &  =-\left\{  \frac{1}{\sqrt{2}}\left(  \left\vert \mathbf{v}%
_{2}\right\rangle +\left\vert \mathbf{v}_{2}^{\ast}\right\rangle \right)
\right\}  ,
\end{align}
i.e.%

\begin{align}
\mathbb{A}\left\vert \mathbf{s}_{1}\right\rangle  &  =\left\vert
\mathbf{s}_{3}\right\rangle \nonumber\\
\mathbb{A}\left\vert \mathbf{s}_{2}\right\rangle  &  =\left\vert
\mathbf{s}_{4}\right\rangle \\
\mathbb{A}\left\vert \mathbf{s}_{3}\right\rangle  &  =-\left\vert
\mathbf{s}_{1}\right\rangle \\
\mathbb{A}\left\vert \mathbf{s}_{4}\right\rangle  &  =-\left\vert
\mathbf{s}_{2}\right\rangle .
\end{align}
$\rfloor$

Thus
\begin{equation}
\mathbb{A=}\sum_{1\leq j\leq\mathbf{s}}\alpha_{j}\left\{  \left\vert
\mathbf{s}_{j}\right\rangle \left\langle \mathbf{s}_{j+\mathbf{s}}\right\vert
-\left\vert \mathbf{s}_{_{j+\mathbf{s}}}\right\rangle \left\langle
\mathbf{s}_{j}\right\vert \right\}  ,
\end{equation}
$\lceil$%
\begin{equation}
\mathbb{A}=\left(  \left\vert \mathbf{s}_{3}\right\rangle \left\langle
\mathbf{s}_{1}\right\vert +\left\vert \mathbf{s}_{4}\right\rangle \left\langle
\mathbf{s}_{2}\right\vert \right)  -\left(  \left\vert \mathbf{s}%
_{1}\right\rangle \left\langle \mathbf{s}_{3}\right\vert +\left\vert
\mathbf{s}_{2}\right\rangle \left\langle \mathbf{s}_{4}\right\vert \right)
\end{equation}
$\rfloor$

which displays the desired result%

\begin{align}
\left\vert g\right\rangle  &  =\frac{1}{2}\sum_{1\leq j_{1},j_{2}\leq
r}A_{j_{1}j_{2}}\left\vert e_{j_{1}}e_{j_{2}}\right\rangle =\frac{1}{2}%
\sum_{1\leq j_{1},j_{2}\leq r}\left[  \mathbb{O}\left(
\begin{array}
[c]{cc}%
\mathbb{O} & \boldsymbol{\alpha}\\
-\boldsymbol{\alpha} & \mathbb{O}%
\end{array}
\right)  \mathbb{O}^{t}\right]  _{j_{1}j_{2}}\left\vert e_{j_{1}}e_{j_{2}%
}\right\rangle \\
&  =\frac{1}{2}\sum_{1\leq j_{1},j_{2}\leq r}\left[  \left(
\begin{array}
[c]{cc}%
\mathbb{O} & \boldsymbol{\alpha}\\
-\boldsymbol{\alpha} & \mathbb{O}%
\end{array}
\right)  \right]  _{j_{1}j_{2}}\left\vert Oe_{j_{1}}Oe_{j_{2}}\right\rangle \\
&  =\sum_{1\leq j\leq s}\alpha_{j}\left\vert Oe_{j}Oe_{j+s}\right\rangle
=\sum_{1\leq j\leq s}\alpha_{j}\left\vert s_{j}s_{j+s}\right\rangle ,
\end{align}
$\lceil$%
\begin{align}
&  \frac{1}{2}\sum_{1\leq j_{1},j_{2}\leq r}\left[  \mathbb{O}\left(
\begin{array}
[c]{cc}%
\mathbb{O} & \boldsymbol{\alpha}\\
-\boldsymbol{\alpha} & \mathbb{O}%
\end{array}
\right)  \mathbb{O}^{t}\right]  _{j_{1}j_{2}}\left\vert e_{j_{1}}e_{j_{2}%
}\right\rangle =\frac{1}{2}\sum_{\substack{1\leq j_{1},j_{2}\leq r\\1\leq
k_{1},k_{2}\leq r}}\left\{  \mathbb{O}_{j_{1}k_{1}}\beta_{k_{1}k_{2}%
}\mathbb{O}_{j_{2}k_{2}}\left\vert e_{j_{1}}e_{j_{2}}\right\rangle \right\} \\
&  =\frac{1}{2}\sum_{1\leq k_{1},k_{2}\leq r}\left\{  \beta_{k_{1}k_{2}%
}\left\vert \left(  \sum_{1\leq j_{1},j_{2}\leq r}e_{j_{1}}\mathbb{O}%
_{j_{1}k_{1}}\right)  \left(  \sum_{1\leq j_{1},j_{2}\leq r}e_{j_{2}%
}\mathbb{O}_{j_{2}k_{2}}\right)  \right\rangle \right\} \\
&  =\frac{1}{2}\sum_{1\leq k\leq s}\left\{  \beta_{kk+s}\left\vert \left(
\sum_{1\leq j_{1}\leq r}e_{j_{1}}\mathbb{O}_{j_{1}k}\right)  \left(
\sum_{1\leq j_{2}\leq r}e_{j_{2}}\mathbb{O}_{j_{2}k+s}\right)  \right\rangle
\right\}  +\frac{1}{2}\sum_{1\leq k\leq s}\left\{  \beta_{k+sk}\left\vert
\left(  \sum_{1\leq j_{1}\leq r}e_{j_{1}}\mathbb{O}_{j_{1}k+s}\right)  \left(
\sum_{1\leq j_{2}\leq r}e_{j_{2}}\mathbb{O}_{j_{2}k}\right)  \right\rangle
\right\} \\
&  =\frac{1}{2}\sum_{1\leq k\leq s}\left\{  \alpha_{k}\left\vert \left(
\sum_{1\leq j_{1}\leq r}e_{j_{1}}\mathbb{O}_{j_{1}k}\right)  \left(
\sum_{1\leq j_{2}\leq r}e_{j_{2}}\mathbb{O}_{j_{2}k+s}\right)  \right\rangle
\right\}  -\frac{1}{2}\sum_{1\leq k\leq s}\left\{  \alpha_{k}\left\vert
\left(  \sum_{1\leq j_{1}\leq r}e_{j_{1}}\mathbb{O}_{j_{1}k+s}\right)  \left(
\sum_{1\leq j_{2}\leq r}e_{j_{2}}\mathbb{O}_{j_{2}k}\right)  \right\rangle
\right\}
\end{align}
$\rfloor$

i.e. the real cannonical form of a real geminal is
\begin{equation}
\left\vert g\right\rangle =\sum\limits_{1\leq j\leq s}\alpha_{j}\left\vert
s_{j}s_{j+s}\right\rangle ,
\end{equation}
while the complex cannonical form of a real geminal is given by
\begin{equation}
\left\vert g\right\rangle =-i\sum_{1\leq j\leq s}\alpha_{j}\left\vert
v_{j}v_{j}^{\ast}\right\rangle .
\end{equation}


$\lceil$%
\begin{align}
\left\vert g\right\rangle  &  =\frac{1}{2}\sum_{1\leq j_{1},j_{2}\leq
r}A_{j_{1}j_{2}}\left\vert e_{j_{1}}e_{j_{2}}\right\rangle =\frac{1}{2}%
\sum_{1\leq j_{1},j_{2}\leq r}\left[  -i\mathbb{V}\left(
\begin{array}
[c]{cc}%
\boldsymbol{\alpha} & \mathbb{O}\\
\mathbb{O} & -\boldsymbol{\alpha}%
\end{array}
\right)  \mathbb{V}^{\dagger}\right]  _{j_{1}j_{2}}\left\vert e_{j_{1}%
}e_{j_{2}}\right\rangle =\frac{1}{2}i\sum_{1\leq j_{1},j_{2}\leq r}\left[
\mathbb{V}\left(
\begin{array}
[c]{cc}%
-\boldsymbol{\alpha} & \mathbb{O}\\
\mathbb{O} & \boldsymbol{\alpha}%
\end{array}
\right)  \mathbb{V}^{\ast t}\right]  _{j_{1}j_{2}}\left\vert e_{j_{1}}%
e_{j_{2}}\right\rangle \\
&  =\frac{1}{2}i\sum_{\substack{1\leq j_{1},j_{2}\leq r\\1\leq k_{1},k_{2}\leq
r}}V_{j_{1}k_{1}}\mathbf{\beta}_{k_{1}k_{2}}V_{j_{2}k_{2}}^{\ast}\left\vert
e_{j_{1}}e_{j_{2}}\right\rangle =\frac{1}{2}i\sum_{1\leq k\leq r}%
\mathbf{\beta}_{kk}\left\vert \left(  \sum_{1\leq j_{1}\leq r}e_{j_{1}%
}V_{j_{1}k}\right)  \left(  \sum_{1\leq j_{2}\leq r}e_{j_{2}}V_{j_{2}k}^{\ast
}\right)  \right\rangle \\
&  =-\frac{1}{2}i\sum_{1\leq k\leq s}\alpha_{k}\left\vert \left(  \sum_{1\leq
j_{1}\leq r}e_{j_{1}}V_{j_{1}k}\right)  \left(  \sum_{1\leq j_{2}\leq
r}e_{j_{2}}V_{j_{2}k}^{\ast}\right)  \right\rangle +\frac{1}{2}i\sum_{s+1\leq
k\leq r}\alpha_{k}\left\vert \left(  \sum_{1\leq j_{1}\leq r}e_{j_{1}}%
V_{j_{1}k}\right)  \left(  \sum_{1\leq j_{2}\leq r}e_{j_{2}}V_{j_{2}k}^{\ast
}\right)  \right\rangle \\
&  =-\frac{1}{2}i\sum_{1\leq k\leq s}\alpha_{k}\left\vert \left(  \sum_{1\leq
j_{1}\leq r}e_{j_{1}}V_{j_{1}k}\right)  \left(  \sum_{1\leq j_{2}\leq
r}e_{j_{2}}V_{j_{2}k}^{\ast}\right)  \right\rangle +\frac{1}{2}i\sum_{1\leq
k\leq s}\alpha_{k}\left\vert \left(  \sum_{1\leq j_{1}\leq r}e_{j_{1}}%
V_{j_{1}k}^{\ast}\right)  \left(  \sum_{1\leq j_{2}\leq r}e_{j_{2}}V_{j_{2}%
k}\right)  \right\rangle \\
&  =-i\sum_{1\leq k\leq s}\alpha_{k}\left\vert \left(  \sum_{1\leq j_{1}\leq
r}e_{j_{1}}V_{j_{1}k}\right)  \left(  \sum_{1\leq j_{2}\leq r}e_{j_{2}%
}V_{j_{2}k}^{\ast}\right)  \right\rangle =-i\sum_{1\leq k\leq s}\alpha
_{k}\left\vert \left(  Ve_{k}\right)  \left(  V^{\ast}e_{k}\right)
\right\rangle =-i\sum_{1\leq k\leq s}\alpha_{k}\left\vert v_{k}v_{k}^{\ast
}\right\rangle .
\end{align}
$\rfloor$

The bases $\left\{  \left\vert \mathbf{s}_{j}\right\rangle ;1\leq j\leq
r\right\}  $ produces the canonical matrix representation of $\mathbb{A}$
\begin{equation}
A\left(  \left\vert \mathbf{s}\right\rangle \right)  =\left\vert
\mathbf{s}\right\rangle \left(
\begin{array}
[c]{cc}%
\mathbb{O} & \boldsymbol{\alpha}\\
-\boldsymbol{\alpha} & \mathbb{O}%
\end{array}
\right)
\end{equation}
and we have the expansions
\begin{align}
\mathbb{A}  &  =\left\vert \mathbf{s}\right\rangle \left(
\begin{array}
[c]{cc}%
\mathbb{O} & \boldsymbol{\alpha}\\
-\boldsymbol{\alpha} & \mathbb{O}%
\end{array}
\right)  \left\langle \mathbf{s}\right\vert =\left\vert \mathbf{s}%
\right\rangle \left(
\begin{array}
[c]{cc}%
\boldsymbol{\alpha}^{\frac{1}{2}} & \mathbb{O}\\
\mathbb{O} & \boldsymbol{\alpha}^{\frac{1}{2}}%
\end{array}
\right)  \left(
\begin{array}
[c]{cc}%
\mathbb{O} & \mathbb{I}_{s}\\
-\mathbb{I}_{s} & \mathbb{O}%
\end{array}
\right)  \left(
\begin{array}
[c]{cc}%
\boldsymbol{\alpha}^{\frac{1}{2}} & \mathbb{O}\\
\mathbb{O} & \boldsymbol{\alpha}^{\frac{1}{2}}%
\end{array}
\right)  \left\langle \mathbf{s}\right\vert \label{reex}\\
&  =\left\vert \mathbf{\sigma}\right\rangle \left(
\begin{array}
[c]{cc}%
\mathbb{O} & \mathbb{I}_{s}\\
-\mathbb{I}_{s} & \mathbb{O}%
\end{array}
\right)  \left\langle \mathbf{\sigma}\right\vert ,
\end{align}
where
\begin{equation}
\left\vert \mathbf{\sigma}\right\rangle =\left\vert \mathbf{s}\right\rangle
\left(
\begin{array}
[c]{cc}%
\boldsymbol{\alpha}^{\frac{1}{2}} & \mathbb{O}\\
\mathbb{O} & \boldsymbol{\alpha}^{\frac{1}{2}}%
\end{array}
\right)
\end{equation}
and
\begin{align}
\left\vert \mathbf{s}\right\rangle  &  =O\left\vert e_{1},\cdots\cdots
,e_{2s}\right\rangle =\left\vert e_{1},\cdots\cdots,e_{2s}\right\rangle
\mathbb{O}\\
\left\langle \left.  \mathbf{\sigma}\right\vert \mathbf{\sigma}\right\rangle
&  =\left(
\begin{array}
[c]{cc}%
\boldsymbol{\alpha} & \mathbb{O}\\
\mathbb{O} & \boldsymbol{\alpha}%
\end{array}
\right)  .
\end{align}
$\lceil$If
\begin{equation}
A\left(  \left\vert \mathbf{s}\right\rangle \right)  =\left\vert
\mathbf{s}\right\rangle \left(
\begin{array}
[c]{cc}%
\mathbb{O} & -\mathbb{I}_{2}\\
\mathbb{I}_{2} & \mathbb{O}%
\end{array}
\right)
\end{equation}
then%
\begin{align}
\mathbb{A}  &  =\left\vert \mathbf{s}_{1},\mathbf{s}_{2},\mathbf{s}%
_{3},\mathbf{s}_{4}\right\rangle \left(
\begin{array}
[c]{cc}%
\mathbb{O} & -\mathbb{I}_{2}\\
\mathbb{I}_{2} & \mathbb{O}%
\end{array}
\right)  \left\langle \mathbf{s}_{1},\mathbf{s}_{2},\mathbf{s}_{3}%
,\mathbf{s}_{4}\right\vert =\left\vert \mathbf{s}_{3},\mathbf{s}%
_{4},-\mathbf{s}_{1},-\mathbf{s}_{2}\right\rangle \left\langle \mathbf{s}%
_{1},\mathbf{s}_{2},\mathbf{s}_{3},\mathbf{s}_{4}\right\vert \\
&  =\left\vert \mathbf{s}_{3}\right\rangle \left\langle \mathbf{s}%
_{1}\right\vert +\left\vert \mathbf{s}_{4}\right\rangle \left\langle
\mathbf{s}_{2}\right\vert -\left\vert \mathbf{s}_{1}\right\rangle \left\langle
\mathbf{s}_{3}\right\vert -\left\vert \mathbf{s}_{2}\right\rangle \left\langle
\mathbf{s}_{4}\right\vert .
\end{align}
$\rfloor$

The row vector, $\left\vert \mathbf{s}_{1},\cdots,\mathbf{s}_{\mathbf{s}%
},\mathbf{s}_{\mathbf{s}+1},\cdots,\mathbf{s}_{\mathbf{s}+\mathbf{s}%
}\right\rangle ,$ of basis vectors, $\left\{  \left\vert \mathbf{s}%
_{1}\right\rangle ,\cdots,\left\vert \mathbf{s}_{\mathbf{s}}\right\rangle
,\left\vert \mathbf{s}_{\mathbf{s}+1}\right\rangle ,\cdots,\left\vert
\mathbf{s}_{\mathbf{s}+\mathbf{s}}\right\rangle \right\}  $, can be viewed as
a real orthogonal matrix, $\mathbb{O},$ whose columns are the real basis
vectors $\left\{  \left\vert \mathbf{s}_{1}\right\rangle ,\cdots\left\vert
\mathbf{s}_{\mathbf{s}}\right\rangle ,\left\vert \mathbf{s}_{\mathbf{s}%
+1}\right\rangle ,\cdots\left\vert \mathbf{s}_{\mathbf{s}+\mathbf{s}%
}\right\rangle \right\}  ,$ such that
\begin{align}
O\left\vert e_{1},\cdots\cdots,e_{2s}\right\rangle  &  =\left\vert
e_{1},\cdots\cdots,e_{2s}\right\rangle \mathbb{O}=\left\vert s_{1}%
,\cdots,s_{s},s_{s+1},\cdots,s_{s+s}\right\rangle \\
\left\vert \mathbf{s}_{1},\cdots,\mathbf{s}_{\mathbf{s}},\mathbf{s}%
_{\mathbf{s}+1},\cdots,\mathbf{s}_{\mathbf{s}+\mathbf{s}}\right\rangle  &
\equiv\mathbb{O}\equiv R_{e}\left(  O\right)
\end{align}
and
\begin{align}
\mathbb{A}  &  =\mathbb{O}\left(
\begin{array}
[c]{cc}%
\mathbb{O} & \boldsymbol{\alpha}\\
-\boldsymbol{\alpha} & \mathbb{O}%
\end{array}
\right)  \mathbb{O}^{t}\\
A\left\vert s_{1},\cdots,s_{s},s_{s+1},\cdots,s_{s+s}\right\rangle  &
=\left\vert s_{1},\cdots,s_{s},s_{s+1},\cdots,s_{s+s}\right\rangle \left(
\begin{array}
[c]{cc}%
\mathbb{O} & \boldsymbol{\alpha}\\
-\boldsymbol{\alpha} & \mathbb{O}%
\end{array}
\right)  .
\end{align}
$\lceil$Then%
\begin{align}
\mathbb{A}  &  =\mathbb{O}\left(
\begin{array}
[c]{cc}%
\boldsymbol{\alpha}^{\frac{1}{2}} & \mathbb{O}\\
\mathbb{O} & \boldsymbol{\alpha}^{\frac{1}{2}}%
\end{array}
\right)  \left(
\begin{array}
[c]{cc}%
\mathbb{O} & \mathbb{I}_{s}\\
-\mathbb{I}_{s} & \mathbb{O}%
\end{array}
\right)  \left(
\begin{array}
[c]{cc}%
\boldsymbol{\alpha}^{\frac{1}{2}} & \mathbb{O}\\
\mathbb{O} & \boldsymbol{\alpha}^{\frac{1}{2}}%
\end{array}
\right)  \mathbb{O}^{t}\\
&  =\mathbb{Q}\left(
\begin{array}
[c]{cc}%
\mathbb{O} & \mathbb{I}_{s}\\
-\mathbb{I}_{s} & \mathbb{O}%
\end{array}
\right)  \mathbb{Q}^{t}%
\end{align}
but%
\begin{align}
\mathbb{QQ}^{t}  &  =\mathbb{O}\left(
\begin{array}
[c]{cc}%
\boldsymbol{\alpha}^{\frac{1}{2}} & \mathbb{O}\\
\mathbb{O} & \boldsymbol{\alpha}^{\frac{1}{2}}%
\end{array}
\right)  \left(
\begin{array}
[c]{cc}%
\boldsymbol{\alpha}^{\frac{1}{2}} & \mathbb{O}\\
\mathbb{O} & \boldsymbol{\alpha}^{\frac{1}{2}}%
\end{array}
\right)  \mathbb{O}^{t}\\
&  =\mathbb{O}\left(
\begin{array}
[c]{cc}%
\boldsymbol{\alpha} & \mathbb{O}\\
\mathbb{O} & \boldsymbol{\alpha}%
\end{array}
\right)  \mathbb{O}^{t}%
\end{align}
and in general
\begin{equation}
\mathbb{O}\left(
\begin{array}
[c]{cc}%
\boldsymbol{\alpha} & \mathbb{O}\\
\mathbb{O} & \boldsymbol{\alpha}%
\end{array}
\right)  \mathbb{O}^{t}\neq\mathbb{I}_{2s}.
\end{equation}
$\rfloor$

$\lceil$Let
\begin{equation}
\mathbb{TAT}^{t}=\mathbb{A}%
\end{equation}
then
\begin{gather}
\mathbb{TO}\left(
\begin{array}
[c]{cc}%
\boldsymbol{\alpha}^{\frac{1}{2}} & \mathbb{O}\\
\mathbb{O} & \boldsymbol{\alpha}^{\frac{1}{2}}%
\end{array}
\right)  \left(
\begin{array}
[c]{cc}%
\mathbb{O} & \mathbb{I}_{s}\\
-\mathbb{I}_{s} & \mathbb{O}%
\end{array}
\right)  \left(
\begin{array}
[c]{cc}%
\boldsymbol{\alpha}^{\frac{1}{2}} & \mathbb{O}\\
\mathbb{O} & \boldsymbol{\alpha}^{\frac{1}{2}}%
\end{array}
\right)  \mathbb{O}^{t}\mathbb{T}^{t}=\mathbb{O}\left(
\begin{array}
[c]{cc}%
\boldsymbol{\alpha}^{\frac{1}{2}} & \mathbb{O}\\
\mathbb{O} & \boldsymbol{\alpha}^{\frac{1}{2}}%
\end{array}
\right)  \left(
\begin{array}
[c]{cc}%
\mathbb{O} & \mathbb{I}_{s}\\
-\mathbb{I}_{s} & \mathbb{O}%
\end{array}
\right)  \left(
\begin{array}
[c]{cc}%
\boldsymbol{\alpha}^{\frac{1}{2}} & \mathbb{O}\\
\mathbb{O} & \boldsymbol{\alpha}^{\frac{1}{2}}%
\end{array}
\right)  \mathbb{O}^{t}\\
\Downarrow\\
\left(
\begin{array}
[c]{cc}%
\boldsymbol{\alpha}^{-\frac{1}{2}} & \mathbb{O}\\
\mathbb{O} & \boldsymbol{\alpha}^{-\frac{1}{2}}%
\end{array}
\right)  \mathbb{O}^{t}\mathbb{TO}\left(
\begin{array}
[c]{cc}%
\boldsymbol{\alpha}^{\frac{1}{2}} & \mathbb{O}\\
\mathbb{O} & \boldsymbol{\alpha}^{\frac{1}{2}}%
\end{array}
\right)  \left(
\begin{array}
[c]{cc}%
\mathbb{O} & \mathbb{I}_{s}\\
-\mathbb{I}_{s} & \mathbb{O}%
\end{array}
\right)  \left(
\begin{array}
[c]{cc}%
\boldsymbol{\alpha}^{\frac{1}{2}} & \mathbb{O}\\
\mathbb{O} & \boldsymbol{\alpha}^{\frac{1}{2}}%
\end{array}
\right)  \mathbb{O}^{t}\mathbb{T}^{t}\mathbb{O}\left(
\begin{array}
[c]{cc}%
\boldsymbol{\alpha}^{-\frac{1}{2}} & \mathbb{O}\\
\mathbb{O} & \boldsymbol{\alpha}^{-\frac{1}{2}}%
\end{array}
\right)  =\left(
\begin{array}
[c]{cc}%
\mathbb{O} & \mathbb{I}_{s}\\
-\mathbb{I}_{s} & \mathbb{O}%
\end{array}
\right) \\
\Downarrow\\
\mathbb{K}\left(
\begin{array}
[c]{cc}%
\mathbb{O} & \mathbb{I}_{s}\\
-\mathbb{I}_{s} & \mathbb{O}%
\end{array}
\right)  \mathbb{K}^{t}=\left(
\begin{array}
[c]{cc}%
\mathbb{O} & \mathbb{I}_{s}\\
-\mathbb{I}_{s} & \mathbb{O}%
\end{array}
\right)
\end{gather}
i.e.
\begin{equation}
\mathbb{K\in}Sp\left(  2s,\mathbb{R}\right)  ,
\end{equation}
where
\begin{equation}
\mathbb{K=}\left(
\begin{array}
[c]{cc}%
\boldsymbol{\alpha}^{-\frac{1}{2}} & \mathbb{O}\\
\mathbb{O} & \boldsymbol{\alpha}^{-\frac{1}{2}}%
\end{array}
\right)  \mathbb{O}^{t}\mathbb{TO}\left(
\begin{array}
[c]{cc}%
\boldsymbol{\alpha}^{\frac{1}{2}} & \mathbb{O}\\
\mathbb{O} & \boldsymbol{\alpha}^{\frac{1}{2}}%
\end{array}
\right)  .
\end{equation}
This shows that%
\begin{equation}
\mathbb{O}\left(
\begin{array}
[c]{cc}%
\boldsymbol{\alpha}^{\frac{1}{2}} & \mathbb{O}\\
\mathbb{O} & \boldsymbol{\alpha}^{\frac{1}{2}}%
\end{array}
\right)  \mathbb{T}\left(
\begin{array}
[c]{cc}%
\boldsymbol{\alpha}^{-\frac{1}{2}} & \mathbb{O}\\
\mathbb{O} & \boldsymbol{\alpha}^{-\frac{1}{2}}%
\end{array}
\right)  \mathbb{O}^{t}=\mathbb{STS}^{-1}=\mathbb{K}%
\end{equation}
so that
\begin{equation}
\mathbb{K\in}Sp\left(  2s,\mathbb{R}\right)
\end{equation}
\emph{but note that in general}
\begin{equation}
\mathbb{K\notin}O\left(  2s,\mathbb{R}\right)  .
\end{equation}
$\rfloor$

$\lceil$If
\begin{equation}
\mathbb{A}=\left(
\begin{array}
[c]{cc}%
\mathbb{O} & -\mathbb{I}_{2}\\
\mathbb{I}_{2} & \mathbb{O}%
\end{array}
\right)
\end{equation}
then
\begin{align}
\mathbb{A}  &  =\left(
\begin{array}
[c]{cc}%
\mathbb{O} & -\mathbb{I}_{2}\\
\mathbb{I}_{2} & \mathbb{O}%
\end{array}
\right)  =\mathbb{O}\left(
\begin{array}
[c]{cc}%
\mathbb{O} & -\mathbb{I}_{2}\\
\mathbb{I}_{2} & \mathbb{O}%
\end{array}
\right)  \mathbb{O}^{t}\\
&  \Longleftrightarrow O\in Sp\left(  4,\mathbb{R}\right)  \cap O\left(
4,\mathbb{R}\right)  .
\end{align}
$\rfloor$

The complex basis $\left\{  \left\vert \mathbf{v}_{j}\right\rangle ,\left\vert
\mathbf{v}_{j}^{\ast}\right\rangle ;1\leq j\leq r\right\}  $ diagonalizes
$\mathbb{A}$ to give
\begin{equation}
\mathbb{A}\left(  \left\vert \mathbf{v}\right\rangle ,\left\vert
\mathbf{v}^{\ast}\right\rangle \right)  =\left(  \left\vert \mathbf{v}%
\right\rangle ,\left\vert \mathbf{v}^{\ast}\right\rangle \right)  \left(
\begin{array}
[c]{cc}%
-i\boldsymbol{\alpha} & \mathbb{O}\\
\mathbb{O} & i\boldsymbol{\alpha}%
\end{array}
\right)
\end{equation}
and produces the expansion%
\begin{equation}
\mathbb{A}=\left(  \left\vert \mathbf{v}\right\rangle ,\left\vert
\mathbf{v}^{\ast}\right\rangle \right)  \left(
\begin{array}
[c]{cc}%
-i\boldsymbol{\alpha} & \mathbb{O}\\
\mathbb{O} & i\boldsymbol{\alpha}%
\end{array}
\right)  \left(
\begin{array}
[c]{c}%
\left\langle \mathbf{v}\right\vert \\
\left\langle \mathbf{v}^{\ast}\right\vert
\end{array}
\right)  =\mathbb{V}\left(
\begin{array}
[c]{cc}%
-i\boldsymbol{\alpha} & \mathbb{O}\\
\mathbb{O} & i\boldsymbol{\alpha}%
\end{array}
\right)  \mathbb{V}^{\dagger}. \label{ccex}%
\end{equation}
$\lceil$%
\begin{equation}
\mathbb{A}\left(  \left\vert \mathbf{v}\right\rangle ,\left\vert
\mathbf{v}^{\ast}\right\rangle \right)  =\left(  \left\vert \mathbf{v}%
\right\rangle ,\left\vert \mathbf{v}^{\ast}\right\rangle \right)  \left(
\begin{array}
[c]{cc}%
i\mathbb{I}_{2} & \mathbb{O}\\
\mathbb{O} & -i\mathbb{I}_{2}%
\end{array}
\right)
\end{equation}%
\begin{align}
\mathbb{A}  &  =\left(  \left\vert \mathbf{v}_{1},\mathbf{v}_{2}%
,\mathbf{v}_{1}^{\ast},\mathbf{v}_{2}^{\ast}\right\rangle \right)  \left(
\begin{array}
[c]{cc}%
i\mathbb{I}_{2} & \mathbb{O}\\
\mathbb{O} & -i\mathbb{I}_{2}%
\end{array}
\right)  \left(  \left\langle \mathbf{v}_{1},\mathbf{v}_{2},\mathbf{v}%
_{1}^{\ast},\mathbf{v}_{2}^{\ast}\right\vert \right)  =\left(  \left\vert
i\mathbf{v}_{1},i\mathbf{v}_{2},-i\mathbf{v}_{1}^{\ast},-i\mathbf{v}_{2}%
^{\ast}\right\rangle \right)  \left(  \left\langle \mathbf{v}_{1}%
,\mathbf{v}_{2},\mathbf{v}_{1}^{\ast},\mathbf{v}_{2}^{\ast}\right\vert \right)
\\
&  =\left(  \left\vert i\mathbf{v}_{1},i\mathbf{v}_{2},-i\mathbf{v}_{1}^{\ast
},-i\mathbf{v}_{2}^{\ast}\right\rangle \right)  \left(  \left\vert
i\mathbf{v}_{1},i\mathbf{v}_{2},-i\mathbf{v}_{1}^{\ast},-i\mathbf{v}_{2}%
^{\ast}\right\rangle \right)  ^{\dagger}\\
&  =i\left\{  \left\vert \mathbf{v}_{1}\right\rangle \left\langle
\mathbf{v}_{1}\right\vert +\left\vert \mathbf{v}_{2}\right\rangle \left\langle
\mathbf{v}_{2}\right\vert -\left\vert \mathbf{v}_{1}^{\ast}\right\rangle
\left\langle \mathbf{v}_{1}^{\ast}\right\vert -\left\vert \mathbf{v}_{2}%
^{\ast}\right\rangle \left\langle \mathbf{v}_{2}^{\ast}\right\vert \right\}  .
\end{align}
$\rfloor$

The row vector of basis vectors, $\left\vert \mathbf{v}_{1},\cdots
,\mathbf{v}_{\mathbf{s}},\mathbf{v}_{1}^{\ast},\cdots,\mathbf{v}_{\mathbf{s}%
}^{\ast}\right\rangle $, can be viewed as a unitary matrix whose columns are
the basis vectors $\left\{  \mathbf{v}_{1},\cdots,\mathbf{v}_{\mathbf{s}%
},\mathbf{v}_{1}^{\ast},\cdots,\mathbf{v}_{\mathbf{s}}^{\ast}\right\}  $. such
that
\begin{align}
V\left\vert e_{1},\cdots\cdots,e_{2\mathbf{s}}\right\rangle  &  =\left\vert
e_{1},\cdots\cdots,e_{2\mathbf{s}}\right\rangle \mathbb{V}=\left\vert
v_{1},\cdots,v_{\mathbf{s}},v_{1}^{\ast},\cdots,v_{\mathbf{s}}^{\ast
}\right\rangle \\
\mathbb{V}  &  =\left(
\begin{array}
[c]{cc}%
\mathbb{V}_{s} & \mathbb{V}_{s}^{\ast}%
\end{array}
\right)  =\left\vert \mathbf{v}_{1},\cdots,\mathbf{v}_{\mathbf{s}}%
,\mathbf{v}_{1}^{\ast},\cdots,\mathbf{v}_{\mathbf{s}}^{\ast}\right\rangle
\end{align}
and
\begin{align}
\mathbb{A}  &  =i\mathbb{V}\left(
\begin{array}
[c]{cc}%
-\boldsymbol{\alpha} & \mathbb{O}\\
\mathbb{O} & \boldsymbol{\alpha}%
\end{array}
\right)  \mathbb{V}^{\dagger}\\
\mathbb{A}  &  =i\left\vert \mathbf{v}_{1},\cdots,\mathbf{v}_{\mathbf{s}%
},\mathbf{v}_{1}^{\ast},\cdots,\mathbf{v}_{\mathbf{s}}^{\ast}\right\rangle
\left(
\begin{array}
[c]{cc}%
-\boldsymbol{\alpha} & \mathbb{O}\\
\mathbb{O} & \boldsymbol{\alpha}%
\end{array}
\right)  \left\langle \mathbf{v}_{1},\cdots,\mathbf{v}_{\mathbf{s}}%
,\mathbf{v}_{1}^{\ast},\cdots,\mathbf{v}_{\mathbf{s}}^{\ast}\right\vert .
\end{align}
This gives
\begin{equation}
\mathbb{A}=i\left\vert \mathbf{\upsilon}_{1},\cdots,\mathbf{\upsilon
}_{\mathbf{s}},\mathbf{\upsilon}_{1}^{\ast},\cdots,\mathbf{\upsilon
}_{\mathbf{s}}^{\ast}\right\rangle \left(
\begin{array}
[c]{cc}%
\mathbb{I}_{s} & \mathbb{O}\\
\mathbb{O} & -\mathbb{I}_{s}%
\end{array}
\right)  \left\langle \mathbf{\upsilon}_{1},\cdots,\mathbf{\upsilon
}_{\mathbf{s}},\mathbf{\upsilon}_{1}^{\ast},\cdots,\mathbf{\upsilon
}_{\mathbf{s}}^{\ast}\right\vert ,
\end{equation}
where
\begin{equation}
\left\vert \mathbf{\upsilon}_{1},\cdots,\mathbf{\upsilon}_{\mathbf{s}%
},\mathbf{\upsilon}_{1}^{\ast},\cdots,\mathbf{\upsilon}_{\mathbf{s}}^{\ast
}\right\rangle =\left\vert \mathbf{v}_{1},\cdots,\mathbf{v}_{\mathbf{s}%
},\mathbf{v}_{1}^{\ast},\cdots,\mathbf{v}_{\mathbf{s}}^{\ast}\right\rangle
\left(
\begin{array}
[c]{cc}%
\boldsymbol{\alpha}^{\frac{1}{2}} & \mathbb{O}\\
\mathbb{O} & \boldsymbol{\alpha}^{\frac{1}{2}}%
\end{array}
\right)  .
\end{equation}
Note that the expansions of eqs. $\left(  \ref{reex}\right)  $ and $\left(
\ref{ccex}\right)  $ both give $\mathbb{A}.$ Thus in eq. $\left(
\ref{reex}\right)  $ all the imaginary terms cancel.

\subsection{Resolutions of the identity}

The following are resolutions of the identity%
\begin{equation}
I=\left\vert \mathbf{e}\right)  \left(  \left.  \mathbf{e}\right\vert
\mathbf{e}\right)  ^{-1}\left(  \mathbf{e}\right\vert =\left\vert
\mathbf{f}\right)  \left(  \left.  \mathbf{f}\right\vert \mathbf{f}\right)
^{-1}\left(  \mathbf{f}\right\vert
\end{equation}%
\begin{equation}
I=\left\vert \mathbf{e}\right)  \left(  \left.  \mathbf{e}\mathcal{K}%
\right\vert \mathbf{e}\right)  ^{-1}\left(  \mathcal{K}^{t}\mathbf{e}%
\right\vert =\left\vert \mathbf{e}\right)  \left(  \left.  \mathbf{e}%
\right\vert \mathbf{e}\right)  _{\mathcal{K}}^{-1}\left(  \mathbf{e}%
\right\vert _{\mathcal{K}},
\end{equation}
which can be seen by observing
\begin{align}
\left\vert x\right)   &  =\left\vert \mathbf{e}\right)  \left(  \left.
\mathbf{e}\right\vert \mathbf{e}\right)  ^{-1}\left(  \mathbf{e}\right\vert
\left\vert x\right)  =\left\vert \mathbf{e}\right)  \left(  \left.
\mathbf{e}\right\vert \mathbf{e}\right)  ^{-1}\left(  \mathbf{e}\right\vert
\left\vert \mathbf{e}\right)  R_{e}\left(  x\right)  =\left\vert
\mathbf{e}\right)  R_{e}\left(  x\right)  =\left\vert x\right) \\
\left\vert x\right)   &  =\left\vert \mathbf{e}\right)  \left(  \left.
\mathbf{e}\mathcal{K}\right\vert \mathbf{e}\right)  ^{-1}\left(
\mathcal{K}^{t}\mathbf{e}\right\vert \left\vert x\right)  =\left\vert
\mathbf{e}\right)  \left(  \left.  \mathbf{e}\mathcal{K}\right\vert
\mathbf{e}\right)  ^{-1}\left(  \mathcal{K}^{t}\mathbf{e}\right\vert
\left\vert \mathbf{e}\right)  \left(  \left.  \mathbf{e}\right\vert
\mathbf{e}\right)  ^{-1}\left(  \mathbf{e}\right\vert \left\vert x\right) \\
&  =\left\vert \mathbf{e}\right)  \left(  \left.  \mathbf{e}\right\vert
\mathbf{e}\right)  ^{-1}\left(  \mathbf{e}\right\vert \left\vert x\right)
=\left\vert \mathbf{e}\right)  \left(  \left.  \mathbf{e}\right\vert
\mathbf{e}\right)  ^{-1}\left(  \mathbf{e}\right\vert \left\vert
\mathbf{e}\right)  R_{e}\left(  x\right)  =\left\vert \mathbf{e}\right)
R_{e}\left(  x\right)  =\left\vert x\right)  .
\end{align}
Using a resolution of the identity one obtains
\begin{equation}
X\left\vert \mathbf{e}\right)  =\left\vert \mathbf{e}\right)  R_{e}\left(
X\right)  =\left\vert \mathbf{e}\right)  \left(  \left.  \mathbf{e}\right\vert
\mathbf{e}\right)  ^{-1}\left(  \left.  \mathbf{e}\right\vert X\mathbf{e}%
\right)
\end{equation}
i.e.
\begin{equation}
R_{e}\left(  X\right)  =\left(  \left.  \mathbf{e}\right\vert \mathbf{e}%
\right)  ^{-1}\left(  \left.  \mathbf{e}\right\vert X\mathbf{e}\right)
=\mathbb{S}_{e}^{-1}M_{e}\left(  X\right)
\end{equation}
and hence
\begin{equation}
\left(  \left.  \mathbf{e}\right\vert X^{-1}\mathbf{e}\right)  =\left(
\left.  \mathbf{e}\right\vert \mathbf{e}\right)  \left(  \left.
\mathbf{e}\right\vert X\mathbf{e}\right)  ^{-1}\left(  \left.  \mathbf{e}%
\right\vert \mathbf{e}\right)
\end{equation}
i.e.
\begin{equation}
M_{e}\left(  X^{-1}\right)  =\mathbb{S}_{e}M_{e}\left(  X\right)
^{-1}\mathbb{S}_{e}%
\end{equation}
$\lceil$%

\begin{equation}
X^{-1}X\left\vert \mathbf{e}\right)  =\left\vert \mathbf{e}\right)
R_{e}\left(  X^{-1}\right)  R_{e}\left(  X\right)  =\left\vert \mathbf{e}%
\right)  \left(  \left.  \mathbf{e}\right\vert \mathbf{e}\right)  ^{-1}\left(
\left.  \mathbf{e}\right\vert X^{-1}\mathbf{e}\right)  \left(  \left.
\mathbf{e}\right\vert \mathbf{e}\right)  ^{-1}\left(  \left.  \mathbf{e}%
\right\vert X\mathbf{e}\right)  =\left\vert \mathbf{e}\right)
\end{equation}
$\rfloor$

One should note the important theorem

\begin{theorem}%
\begin{equation}
R_{e}\left(  X\right)  =\left(  \left.  \mathbf{e}\right\vert \mathbf{e}%
\right)  _{\mathcal{K}}^{-1}\left(  \left.  \mathbf{e}\right\vert
X\mathbf{e}\right)  _{\mathcal{K}}=\mathbb{S}_{\mathcal{K}e}^{-1}%
M_{\mathcal{K}e}\left(  X\right)  =\left(  \left.  \mathbf{e}\right\vert
\mathcal{K}\mathbf{e}\right)  ^{-1}\left(  \left.  \mathbf{e}\right\vert
\mathcal{K}X\mathbf{e}\right)  =\left(  \left.  \mathbf{e}\right\vert
\mathbf{e}\right)  ^{-1}\left(  \left.  \mathbf{e}\right\vert X\mathbf{e}%
\right)
\end{equation}
for any basis $\left\{  e_{1},\cdots\cdots,e_{2s}\right\}  $ and invertible
operator $\mathcal{K}$
\end{theorem}

\begin{proof}%
\begin{equation}
\left(  \left.  \mathbf{e}\right\vert \mathcal{K}\mathbf{e}\right)
^{-1}\left(  \left.  \mathbf{e}\right\vert \mathcal{K}X\mathbf{e}\right)
=\left(  \left.  \mathbf{e}\right\vert \mathcal{K}\mathbf{e}\right)
^{-1}\left(  \left.  \mathbf{e}\right\vert \mathcal{K}\mathbf{e}\right)
\left(  \left.  \mathbf{e}\right\vert \mathbf{e}\right)  ^{-1}\left(  \left.
\mathbf{e}\right\vert X\mathbf{e}\right)  =\left(  \left.  \mathbf{e}%
\right\vert \mathbf{e}\right)  ^{-1}\left(  \left.  \mathbf{e}\right\vert
X\mathbf{e}\right)  .
\end{equation}

\end{proof}

\subsection{Inner Products}

\qquad One can construct symmetric and antisymmetric inner prodcuts of
different signitures by considering non singular symmetric and antisymmetric
operators $\mathcal{K}$ by%
\begin{equation}
\left(  \left.  x\right\vert y\right)  _{\mathcal{K}}=\left(  \left.
x\right\vert \mathcal{K}y\right)  .
\end{equation}
These inner products induce linear functionals
\begin{equation}
\left(  x\right\vert _{\mathcal{K}}=\left(  x\right\vert \mathcal{K}=\left(
\mathcal{K}^{t}x\right\vert ,
\end{equation}
Transposes defined by the different inner products are related by
\begin{align}
\left(  \left.  M^{t}x\right\vert y\right)   &  =\left(  \left.  x\right\vert
My\right)  _{\mathcal{K}}=\left(  \left.  x\right\vert \mathcal{K}My\right)
=\left(  \left.  M^{t}\mathcal{K}^{t}x\right\vert y\right) \\
&  =\left(  \left.  \mathcal{K}^{-t}M^{t}\mathcal{K}^{t}x\right\vert
\mathcal{K}y\right)  =\left(  \left.  \mathcal{K}^{-t}M^{t}\mathcal{K}%
^{t}x\right\vert y\right)  _{\mathcal{K}}%
\end{align}
therefore
\begin{equation}
M^{t_{\mathcal{K}}}=\mathcal{K}^{-t}M^{t}\mathcal{K}^{t}%
\end{equation}
and in particular
\begin{equation}
\mathcal{K}^{t_{\mathcal{K}}}=\mathcal{K}^{-t}\mathcal{K}^{t}\mathcal{K}%
^{t}=\mathcal{K}^{t}.
\end{equation}


\subsection{Representations}

The representation of an operator $T$ wrt the bases $\left\{  e_{1}%
,\cdots\cdots,e_{2s}\right\}  $ is given by
\begin{align}
T\left(  e_{1},\cdots\cdots,e_{2s}\right)   &  =\left(  e_{1},\cdots
\cdots,e_{2s}\right)  R_{e}\left(  T\right) \\
R_{e}\left(  T\right)   &  =\mathbb{S}_{\mathcal{K}e}^{-1}M_{\mathcal{K}%
e}\left(  T\right) \\
M_{\mathcal{K}e}\left(  T\right)   &  =\left(  e_{1},\cdots\cdots
,e_{2s}\left\vert Te_{1},\cdots\cdots,e_{2s}\right.  \right)  _{\mathcal{K}%
}=\left(  e_{1},\cdots\cdots,e_{2s}\left\vert \mathcal{K}Te_{1},\cdots
\cdots,e_{2s}\right.  \right) \\
\mathbb{S}_{\mathcal{K}e}^{-1}  &  =\left(  e_{1},\cdots\cdots,e_{2s}%
\left\vert e_{1},\cdots\cdots,e_{2s}\right.  \right)  _{\mathcal{K}}%
^{-1}=\left(  e_{1},\cdots\cdots,e_{2s}\left\vert \mathcal{K}e_{1}%
,\cdots\cdots,e_{2s}\right.  \right)  ^{-1}.
\end{align}


The representations of the transformation $\mathcal{T}_{fe}$ between the bases
$\left\{  e_{1},\cdots\cdots,e_{2s}\right\}  $ and $\left\{  f_{1}%
,\cdots\cdots,f_{2s}\right\}  $ is given by
\begin{align}
\mathcal{T}_{fe}\left(  e_{1},\cdots\cdots,e_{2s}\right)   &  =\left(
e_{1},\cdots\cdots,e_{2s}\right)  R_{e}\left(  \mathcal{T}_{fe}\right)
=\left(  f_{1},\cdots\cdots,f_{2s}\right) \\
R_{e}\left(  \mathcal{T}_{fe}\right)   &  =\mathbb{S}_{\mathcal{K}e}%
^{-1}M_{\mathcal{K}e}\left(  \mathcal{T}_{fe}\right) \\
M_{\mathcal{K}e}\left(  \mathcal{T}_{fe}\right)   &  =\left(  e_{1}%
,\cdots\cdots,e_{2s}\left\vert \mathcal{T}_{fe}e_{1},\cdots\cdots
,e_{2s}\right.  \right)  _{\mathcal{K}}.
\end{align}


The representations of an operator $T$ wrt to the two different bases is given
by
\begin{align}
T\left(  e_{1},\cdots\cdots,e_{2s}\right)   &  =\left(  e_{1},\cdots
\cdots,e_{2s}\right)  R_{e}\left(  T\right) \\
R_{e}\left(  T\right)   &  =\mathbb{S}_{\mathcal{K}e}^{-1}M_{\mathcal{K}%
e}\left(  T\right) \\
M_{\mathcal{K}e}\left(  T\right)   &  =\left(  e_{1},\cdots\cdots
,e_{2s}\left\vert Te_{1},\cdots\cdots,e_{2s}\right.  \right)  _{\mathcal{K}}%
\end{align}
and%
\begin{align}
T\left(  f_{1},\cdots\cdots,f_{2s}\right)   &  =\left(  f_{1},\cdots
\cdots,f_{2s}\right)  R_{f}\left(  T\right) \\
R_{f}\left(  T\right)   &  =\mathbb{S}_{\mathcal{K}f}^{-1}M_{\mathcal{K}%
}\left(  T\right) \\
M_{\mathcal{K}f}\left(  T\right)   &  =\left(  f_{1},\cdots\cdots
,f_{2s}\left\vert Tf_{1},\cdots\cdots,f_{2s}\right.  \right)  _{\mathcal{K}}.
\end{align}
Hence
\begin{align}
\mathcal{T}_{fe}^{-1}T\left(  f_{1},\cdots\cdots,f_{2s}\right)   &
=\mathcal{T}_{fe}^{-1}\left(  f_{1},\cdots\cdots,f_{2s}\right)  R_{f}\left(
T\right)  =\left(  e_{1},\cdots\cdots,e_{2s}\right)  R_{f}\left(  T\right) \\
&  =\mathcal{T}_{fe}^{-1}T\mathcal{T}_{fe}\left(  e_{1},\cdots\cdots
,e_{2s}\right)  =\left(  e_{1},\cdots\cdots,e_{2s}\right)  R_{e}\left(
\mathcal{T}_{fe}^{-1}\right)  R_{e}\left(  T\right)  R_{e}\left(
\mathcal{T}_{fe}\right)
\end{align}
showing that
\begin{equation}
R_{f}\left(  T\right)  =R_{e}\left(  \mathcal{T}_{fe}^{-1}\right)
R_{e}\left(  T\right)  R_{e}\left(  \mathcal{T}_{fe}\right)  =R_{e}\left(
\mathcal{T}_{fe}\right)  ^{-1}R_{e}\left(  T\right)  R_{e}\left(
\mathcal{T}_{fe}\right)  .
\end{equation}


The relationship between matrices wrt the different inner products and bases
are given by
\begin{align}
&  R_{f}\left(  T\right)  =R_{e}\left(  \mathcal{T}_{fe}^{-1}\right)
R_{e}\left(  T\right)  R_{e}\left(  \mathcal{T}_{fe}\right) \\
&  \mathbb{S}_{\mathcal{K}_{2}f}^{-1}M_{\mathcal{K}_{2}f}\left(  T\right)
=\mathbb{S}_{\mathcal{K}_{1}e}^{-1}M_{\mathcal{K}_{1}e}\left(  \mathcal{T}%
_{fe}^{-1}\right)  \mathbb{S}_{\mathcal{K}_{1}e}^{-1}M_{\mathcal{K}_{1}%
e}\left(  T\right)  \mathbb{S}_{\mathcal{K}_{1}e}^{-1}M_{\mathcal{K}_{1}%
e}\left(  \mathcal{T}_{fe}\right) \\
&  M_{\mathcal{K}_{2}f}\left(  T\right)  =\mathbb{S}_{\mathcal{K}_{2}%
f}\mathbb{S}_{\mathcal{K}_{1}e}^{-1}M_{\mathcal{K}_{1}e}\left(  \mathcal{T}%
_{fe}^{-1}\right)  \mathbb{S}_{\mathcal{K}_{1}e}^{-1}M_{\mathcal{K}_{1}%
e}\left(  T\right)  \mathbb{S}_{\mathcal{K}_{1}e}^{-1}M_{\mathcal{K}_{1}%
e}\left(  \mathcal{T}_{fe}\right) \\
&  \left(  f_{1},\cdots\cdots,f_{2s}\left\vert \mathcal{K}_{2}Tf_{1}%
,\cdots\cdots,f_{2s}\right.  \right)  =\left(  f_{1},\cdots\cdots
,f_{2s}\left\vert \mathcal{K}_{2}f_{1},\cdots\cdots,f_{2s}\right.  \right)
\left(  e_{1},\cdots\cdots,e_{2s}\left\vert \mathcal{K}_{1}e_{1},\cdots
\cdots,e_{2s}\right.  \right)  \left(  e_{1},\cdots\cdots,e_{2s}\left\vert
\mathcal{K}_{1}\mathcal{T}_{fe}^{-1}e_{1},\cdots\cdots,e_{2s}\right.  \right)
\\
&  \times\left(  e_{1},\cdots\cdots,e_{2s}\left\vert \mathcal{K}_{1}%
e_{1},\cdots\cdots,e_{2s}\right.  \right)  ^{-1}\left(  e_{1},\cdots
\cdots,e_{2s}\left\vert \mathcal{K}_{1}Te_{1},\cdots\cdots,e_{2s}\right.
\right)  \left(  e_{1},\cdots\cdots,e_{2s}\left\vert \mathcal{K}_{1}%
e_{1},\cdots\cdots,e_{2s}\right.  \right)  ^{-1}\left(  e_{1},\cdots
\cdots,e_{2s}\left\vert \mathcal{K}_{1}\mathcal{T}_{fe}e_{1},\cdots
\cdots,e_{2s}\right.  \right)  .
\end{align}


Operators can be expanded in the following manner
\begin{equation}
X=\sum_{1\leq j,k\leq r}\left(  \mathbb{S}_{\mathcal{K}_{1}e}^{-1}%
M_{\mathcal{K}_{1}e}\left(  X\right)  \mathbb{S}_{\mathcal{K}_{1}e}%
^{-1}\right)  _{jk}\left\vert e_{j}\right)  \left(  e_{k}\right\vert
_{\mathcal{K}_{1}}=\sum_{1\leq j,k\leq r}\left(  \mathbb{S}_{\mathcal{K}_{1}%
e}^{-1}M_{\mathcal{K}_{1}e}\left(  X\right)  \mathbb{S}_{\mathcal{K}_{1}%
e}^{-1}\right)  _{jk}\left\vert e_{j}\right)  \left(  \mathcal{K}_{1}^{t}%
e_{k}\right\vert .
\end{equation}
$\lceil$%
\begin{align}
\left(  \mathcal{K}_{1}^{t}e_{m}\right\vert X\left\vert e_{l}\right)   &
=\sum_{1\leq j,k\leq r}\left(  \mathcal{K}_{1}^{t}e_{m}\right\vert \left(
\mathbb{S}_{\mathcal{K}_{1}e}^{-1}M_{\mathcal{K}_{1}e}\left(  X\right)
\mathbb{S}_{\mathcal{K}_{1}e}^{-1}\right)  _{jk}\left\vert e_{j}\right)
\left(  \mathcal{K}_{1}^{t}e_{k}\right\vert \left\vert e_{l}\right) \\
&  =\sum_{1\leq j,k\leq r}\left(  \mathbb{S}_{\mathcal{K}_{1}e}^{-1}%
M_{\mathcal{K}_{1}e}\left(  X\right)  \mathbb{S}_{\mathcal{K}_{1}e}%
^{-1}\right)  _{jk}\left(  \mathcal{K}_{1}^{t}e_{m}\right\vert \left\vert
e_{j}\right)  \left(  \mathcal{K}_{1}^{t}e_{k}\right\vert \left\vert
e_{l}\right) \\
&  =\sum_{1\leq j,k\leq r}\left(  \mathcal{K}_{1}^{t}e_{m}\right\vert
\left\vert e_{j}\right)  \left(  \mathbb{S}_{\mathcal{K}_{1}e}^{-1}%
M_{\mathcal{K}_{1}e}\left(  X\right)  \mathbb{S}_{\mathcal{K}_{1}e}%
^{-1}\right)  _{jk}\left(  \mathcal{K}_{1}^{t}e_{k}\right\vert \left\vert
e_{l}\right) \\
&  =\sum_{1\leq j,k\leq r}\left(  \mathbb{S}_{\mathcal{K}_{1}e}^{-1}\right)
_{mj}\left(  \mathbb{S}_{\mathcal{K}_{1}e}^{-1}M_{\mathcal{K}_{1}e}\left(
X\right)  \mathbb{S}_{\mathcal{K}_{1}e}^{-1}\right)  _{jk}\left(
\mathbb{S}_{\mathcal{K}_{1}e}\right)  _{kl}\\
&  =M_{\mathcal{K}_{1}e}\left(  X\right)  _{ml}.
\end{align}
$\rfloor$

$\lceil$Let%
\begin{equation}
\mathbb{TO}\left(
\begin{array}
[c]{cc}%
\boldsymbol{\alpha} & \mathbb{O}\\
\mathbb{O} & \boldsymbol{\alpha}%
\end{array}
\right)  \mathbb{O}^{t}\mathbb{T}^{t}=\left(
\begin{array}
[c]{cc}%
\boldsymbol{\alpha} & \mathbb{O}\\
\mathbb{O} & \boldsymbol{\alpha}%
\end{array}
\right)
\end{equation}
then
\begin{equation}
\left(
\begin{array}
[c]{cc}%
\boldsymbol{\alpha}^{-\frac{1}{2}} & \mathbb{O}\\
\mathbb{O} & \boldsymbol{\alpha}^{-\frac{1}{2}}%
\end{array}
\right)  \mathbb{O}^{t}\mathbb{TO}\left(
\begin{array}
[c]{cc}%
\boldsymbol{\alpha}^{\frac{1}{2}} & \mathbb{O}\\
\mathbb{O} & \boldsymbol{\alpha}^{\frac{1}{2}}%
\end{array}
\right)  \left(
\begin{array}
[c]{cc}%
\boldsymbol{\alpha}^{-\frac{1}{2}} & \mathbb{O}\\
\mathbb{O} & \boldsymbol{\alpha}^{-\frac{1}{2}}%
\end{array}
\right)  \mathbb{O}^{t}\mathbb{T}^{t}\mathbb{O}\left(
\begin{array}
[c]{cc}%
\boldsymbol{\alpha}^{\frac{1}{2}} & \mathbb{O}\\
\mathbb{O} & \boldsymbol{\alpha}^{\frac{1}{2}}%
\end{array}
\right)  =\left(
\begin{array}
[c]{cc}%
\mathbb{I}_{s} & \mathbb{O}\\
\mathbb{O} & \mathbb{I}_{s}%
\end{array}
\right)  .
\end{equation}
Define
\begin{align}
\mathbb{L}  &  =\left(
\begin{array}
[c]{cc}%
\boldsymbol{\alpha}^{-\frac{1}{2}} & \mathbb{O}\\
\mathbb{O} & \boldsymbol{\alpha}^{-\frac{1}{2}}%
\end{array}
\right)  \mathbb{O}^{t}\mathbb{TO}\left(
\begin{array}
[c]{cc}%
\boldsymbol{\alpha}^{\frac{1}{2}} & \mathbb{O}\\
\mathbb{O} & \boldsymbol{\alpha}^{\frac{1}{2}}%
\end{array}
\right) \\
\mathbb{T}  &  =\left(
\begin{array}
[c]{cc}%
\boldsymbol{\alpha}^{\frac{1}{2}} & \mathbb{O}\\
\mathbb{O} & \boldsymbol{\alpha}^{\frac{1}{2}}%
\end{array}
\right)  \mathbb{OLO}^{t}\left(
\begin{array}
[c]{cc}%
\boldsymbol{\alpha}^{-\frac{1}{2}} & \mathbb{O}\\
\mathbb{O} & \boldsymbol{\alpha}^{-\frac{1}{2}}%
\end{array}
\right)
\end{align}
then
\begin{equation}
\mathbb{LL}^{t}=\left(
\begin{array}
[c]{cc}%
\mathbb{I}_{s} & \mathbb{O}\\
\mathbb{O} & \mathbb{I}_{s}%
\end{array}
\right)
\end{equation}
and
\begin{equation}
\mathbb{L\in}O\left(  2s,\mathbb{R}\right)  .
\end{equation}


\subsubsection{Partitioning}

Hence%
\begin{equation}
\mathbb{L}=\left(
\begin{array}
[c]{cc}%
\boldsymbol{\alpha}^{-\frac{1}{2}} & \mathbb{O}\\
\mathbb{O} & \boldsymbol{\alpha}^{-\frac{1}{2}}%
\end{array}
\right)  \left(
\begin{array}
[c]{cc}%
\mathbb{T}_{11} & \mathbb{T}_{12}\\
\mathbb{T}_{21} & \mathbb{T}_{22}%
\end{array}
\right)  \left(
\begin{array}
[c]{cc}%
\boldsymbol{\alpha}^{\frac{1}{2}} & \mathbb{O}\\
\mathbb{O} & \boldsymbol{\alpha}^{\frac{1}{2}}%
\end{array}
\right)  =\left(
\begin{array}
[c]{cc}%
\boldsymbol{\alpha}^{-\frac{1}{2}}\mathbb{T}_{11}\boldsymbol{\alpha}^{\frac
{1}{2}} & \boldsymbol{\alpha}^{-\frac{1}{2}}\mathbb{T}_{12}\boldsymbol{\alpha
}^{\frac{1}{2}}\\
\boldsymbol{\alpha}^{-\frac{1}{2}}\mathbb{T}_{21}\boldsymbol{\alpha}^{\frac
{1}{2}} & \boldsymbol{\alpha}^{-\frac{1}{2}}\mathbb{T}_{22}\boldsymbol{\alpha
}^{\frac{1}{2}}%
\end{array}
\right)  =\left(
\begin{array}
[c]{cc}%
\mathbb{L}_{11} & \mathbb{L}_{12}\\
\mathbb{L}_{21} & \mathbb{L}_{22}%
\end{array}
\right)
\end{equation}
and%
\begin{equation}
\left(
\begin{array}
[c]{cc}%
\mathbb{L}_{11} & \mathbb{L}_{12}\\
\mathbb{L}_{21} & \mathbb{L}_{22}%
\end{array}
\right)  \left(
\begin{array}
[c]{cc}%
\mathbb{L}_{11}^{t} & \mathbb{L}_{21}^{t}\\
\mathbb{L}_{12}^{t} & \mathbb{L}_{22}^{t}%
\end{array}
\right)  =\left(
\begin{array}
[c]{cc}%
\mathbb{L}_{11}\mathbb{L}_{11}^{t}+\mathbb{L}_{12}\mathbb{L}_{12}^{t} &
\mathbb{L}_{11}\mathbb{L}_{21}^{t}+\mathbb{L}_{12}\mathbb{L}_{22}^{t}\\
\mathbb{L}_{21}\mathbb{L}_{11}^{t}+\mathbb{L}_{22}\mathbb{L}_{12}^{t} &
\mathbb{L}_{21}\mathbb{L}_{21}^{t}+\mathbb{L}_{22}\mathbb{L}_{22}^{t}%
\end{array}
\right)  =\left(
\begin{array}
[c]{cc}%
\mathbb{I}_{s} & \mathbb{O}\\
\mathbb{O} & \mathbb{I}_{s}%
\end{array}
\right)  .
\end{equation}


Also
\begin{align}
\left(
\begin{array}
[c]{cc}%
\boldsymbol{\alpha} & \mathbb{O}\\
\mathbb{O} & \boldsymbol{\alpha}%
\end{array}
\right)   &  =\left(
\begin{array}
[c]{cc}%
\mathbb{T}_{11} & \mathbb{T}_{12}\\
\mathbb{T}_{21} & \mathbb{T}_{22}%
\end{array}
\right)  \left(
\begin{array}
[c]{cc}%
\boldsymbol{\alpha} & \mathbb{O}\\
\mathbb{O} & \boldsymbol{\alpha}%
\end{array}
\right)  \left(
\begin{array}
[c]{cc}%
\mathbb{T}_{11}^{t} & \mathbb{T}_{21}^{t}\\
\mathbb{T}_{12}^{t} & \mathbb{T}_{22}^{t}%
\end{array}
\right) \\
&  =\left(
\begin{array}
[c]{cc}%
\mathbb{T}_{11}\boldsymbol{\alpha} & \mathbb{T}_{12}\boldsymbol{\alpha}\\
\mathbb{T}_{21}\boldsymbol{\alpha} & \mathbb{T}_{22}\boldsymbol{\alpha}%
\end{array}
\right)  \left(
\begin{array}
[c]{cc}%
\mathbb{T}_{11}^{t} & \mathbb{T}_{21}^{t}\\
\mathbb{T}_{12}^{t} & \mathbb{T}_{22}^{t}%
\end{array}
\right)  =\left(
\begin{array}
[c]{cc}%
\mathbb{T}_{11}\boldsymbol{\alpha}\mathbb{T}_{11}^{t}+\mathbb{T}%
_{12}\boldsymbol{\alpha}\mathbb{T}_{12}^{t} & \mathbb{T}_{11}%
\boldsymbol{\alpha}\mathbb{T}_{21}^{t}+\mathbb{T}_{12}\boldsymbol{\alpha
}\mathbb{T}_{22}^{t}\\
\mathbb{T}_{21}\boldsymbol{\alpha}\mathbb{T}_{11}^{t}+\mathbb{T}%
_{22}\boldsymbol{\alpha}\mathbb{T}_{12}^{t} & \mathbb{T}_{21}%
\boldsymbol{\alpha}\mathbb{T}_{21}^{t}+\mathbb{T}_{22}\boldsymbol{\alpha
}\mathbb{T}_{22}^{t}%
\end{array}
\right)
\end{align}
i.e.
\begin{equation}
\left(
\begin{array}
[c]{cc}%
\boldsymbol{\alpha} & \mathbb{O}\\
\mathbb{O} & \boldsymbol{\alpha}%
\end{array}
\right)  =\left(
\begin{array}
[c]{cc}%
\mathbb{T}_{11}\boldsymbol{\alpha}\mathbb{T}_{11}^{t}+\mathbb{T}%
_{12}\boldsymbol{\alpha}\mathbb{T}_{12}^{t} & \mathbb{T}_{11}%
\boldsymbol{\alpha}\mathbb{T}_{21}^{t}+\mathbb{T}_{12}\boldsymbol{\alpha
}\mathbb{T}_{22}^{t}\\
\mathbb{T}_{21}\boldsymbol{\alpha}\mathbb{T}_{11}^{t}+\mathbb{T}%
_{22}\boldsymbol{\alpha}\mathbb{T}_{12}^{t} & \mathbb{T}_{21}%
\boldsymbol{\alpha}\mathbb{T}_{21}^{t}+\mathbb{T}_{22}\boldsymbol{\alpha
}\mathbb{T}_{22}^{t}%
\end{array}
\right)  .
\end{equation}
$\rfloor$

\subsubsection{General Forms}

\paragraph{Symplectic}%

\begin{align}
\mathbb{A}  &  =\mathbb{O}\left(
\begin{array}
[c]{cc}%
\mathbb{O} & \boldsymbol{\alpha}\\
-\boldsymbol{\alpha} & \mathbb{O}%
\end{array}
\right)  \mathbb{O}^{t}=\left\vert \mathbf{s}_{1},\cdots,\mathbf{s}%
_{\mathbf{s}},\mathbf{s}_{\mathbf{s}+1},\cdots,\mathbf{s}_{\mathbf{s}%
+\mathbf{s}}\right\rangle \left(
\begin{array}
[c]{cc}%
\mathbb{O} & \boldsymbol{\alpha}\\
-\boldsymbol{\alpha} & \mathbb{O}%
\end{array}
\right)  \left\langle \mathbf{s}_{1},\cdots,\mathbf{s}_{\mathbf{s}}%
,\mathbf{s}_{\mathbf{s}+1},\cdots,\mathbf{s}_{\mathbf{s}+\mathbf{s}%
}\right\vert \\
&  =\left\vert \mathbf{s}_{1},\cdots,\mathbf{s}_{\mathbf{s}},\mathbf{s}%
_{\mathbf{s}+1},\cdots,\mathbf{s}_{\mathbf{s}+\mathbf{s}}\right\rangle \left(
%
\begin{array}
[c]{cc}%
\mathbb{O} & \boldsymbol{\alpha}\\
-\boldsymbol{\alpha} & \mathbb{O}%
\end{array}
\right)  \left\langle \mathbf{s}_{1},\cdots,\mathbf{s}_{\mathbf{s}}%
,\mathbf{s}_{\mathbf{s}+1},\cdots,\mathbf{s}_{\mathbf{s}+\mathbf{s}%
}\right\vert \\
&  =\frac{1}{2}\left\vert \mathbf{v}_{1},\cdots,\mathbf{v}_{\mathbf{s}%
},\mathbf{v}_{1}^{\ast},\cdots,\mathbf{v}_{\mathbf{s}}^{\ast}\right\rangle
\left(
\begin{array}
[c]{cc}%
\mathbb{I}_{s} & -i\mathbb{I}_{s}\\
\mathbb{I}_{s} & i\mathbb{I}_{s}%
\end{array}
\right)  \left(
\begin{array}
[c]{cc}%
\mathbb{O} & \boldsymbol{\alpha}\\
-\boldsymbol{\alpha} & \mathbb{O}%
\end{array}
\right)  \left(
\begin{array}
[c]{cc}%
\mathbb{I}_{s} & \mathbb{I}_{s}\\
i\mathbb{I}_{s} & -i\mathbb{I}_{s}%
\end{array}
\right)  \left\langle \mathbf{v}_{1},\cdots,\mathbf{v}_{\mathbf{s}}%
,\mathbf{v}_{1}^{\ast},\cdots,\mathbf{v}_{\mathbf{s}}^{\ast}\right\vert \\
&  =\left\vert \mathbf{v}_{1},\cdots,\mathbf{v}_{\mathbf{s}},\mathbf{v}%
_{1}^{\ast},\cdots,\mathbf{v}_{\mathbf{s}}^{\ast}\right\rangle i\left(
\begin{array}
[c]{cc}%
-\boldsymbol{\alpha} & \mathbb{O}\\
\mathbb{O} & \boldsymbol{\alpha}%
\end{array}
\right)  \left\langle \mathbf{v}_{1},\cdots,\mathbf{v}_{\mathbf{s}}%
,\mathbf{v}_{1}^{\ast},\cdots,\mathbf{v}_{\mathbf{s}}^{\ast}\right\vert .
\end{align}
i.e.
\begin{equation}
\frac{1}{2}\left(
\begin{array}
[c]{cc}%
\mathbb{I}_{s} & -i\mathbb{I}_{s}\\
\mathbb{I}_{s} & i\mathbb{I}_{s}%
\end{array}
\right)  \left(
\begin{array}
[c]{cc}%
\mathbb{O} & \boldsymbol{\alpha}\\
-\boldsymbol{\alpha} & \mathbb{O}%
\end{array}
\right)  \left(
\begin{array}
[c]{cc}%
\mathbb{I}_{s} & \mathbb{I}_{s}\\
i\mathbb{I}_{s} & -i\mathbb{I}_{s}%
\end{array}
\right)  =i\left(
\begin{array}
[c]{cc}%
-\boldsymbol{\alpha} & \mathbb{O}\\
\mathbb{O} & \boldsymbol{\alpha}%
\end{array}
\right)  .
\end{equation}
Hence
\begin{align}
A  &  =\left\vert e_{1},\cdots\cdots,e_{2s}\right\rangle \mathbb{O}\left(
\begin{array}
[c]{cc}%
\mathbb{O} & \boldsymbol{\alpha}\\
-\boldsymbol{\alpha} & \mathbb{O}%
\end{array}
\right)  \mathbb{O}^{t}\left\langle e_{1},\cdots\cdots,e_{2s}\right\vert
=\left\vert e_{1},\cdots\cdots,e_{2s}\right\rangle i\mathbb{V}\left(
\begin{array}
[c]{cc}%
-\boldsymbol{\alpha} & \mathbb{O}\\
\mathbb{O} & \boldsymbol{\alpha}%
\end{array}
\right)  \mathbb{V}^{\dagger}\left\langle e_{1},\cdots\cdots,e_{2s}\right\vert
\\
&  =\left\vert s_{1},\cdots,s_{s},s_{s+1},\cdots,s_{s+s}\right\rangle \left(
\begin{array}
[c]{cc}%
\mathbb{O} & \boldsymbol{\alpha}\\
-\boldsymbol{\alpha} & \mathbb{O}%
\end{array}
\right)  \left\langle s_{1},\cdots,s_{s},s_{s+1},\cdots,s_{s+s}\right\vert
=\left\vert v_{1},\cdots,v_{s},v_{1}^{\ast},\cdots,v_{s}^{\ast}\right\rangle
i\left(
\begin{array}
[c]{cc}%
-\boldsymbol{\alpha} & \mathbb{O}\\
\mathbb{O} & \boldsymbol{\alpha}%
\end{array}
\right)  \left\langle v_{1},\cdots,v_{\mathbf{s}},v_{1}^{\ast},\cdots
,v_{s}^{\ast}\right\vert \\
&  =\left\vert \tilde{s}_{1},\cdots,\tilde{s}_{s},\tilde{s}_{s+1}%
,\cdots,\tilde{s}_{s+s}\right\rangle \left(
\begin{array}
[c]{cc}%
\mathbb{O} & \mathbb{I}_{s}\\
-\mathbb{I}_{s} & \mathbb{O}%
\end{array}
\right)  \left\langle \tilde{s}_{1},\cdots,\tilde{s}_{s},\tilde{s}%
_{s+1},\cdots,\tilde{s}_{s+s}\right\vert =\left\vert \tilde{v}_{1}%
,\cdots,\tilde{v}_{s},\tilde{v}_{1}^{\ast},\cdots,\tilde{v}_{s}^{\ast
}\right\rangle i\left(
\begin{array}
[c]{cc}%
-\mathbb{I}_{s} & \mathbb{O}\\
\mathbb{O} & \mathbb{I}_{s}%
\end{array}
\right)  \left\langle \tilde{v}_{1},\cdots,\tilde{v}_{s},\tilde{v}_{1}^{\ast
},\cdots,\tilde{v}_{s}^{\ast}\right\vert ,
\end{align}
where
\begin{align}
\left\vert \tilde{s}_{1},\cdots,\tilde{s}_{s},\tilde{s}_{s+1},\cdots,\tilde
{s}_{s+s}\right\rangle  &  =\left\vert s_{1},\cdots,s_{s},s_{s+1}%
,\cdots,s_{s+s}\right\rangle \left(
\begin{array}
[c]{cc}%
\mathbb{O} & \boldsymbol{\alpha}^{\frac{1}{2}}\\
-\boldsymbol{\alpha}^{\frac{1}{2}} & \mathbb{O}%
\end{array}
\right) \\
\left\vert \tilde{v}_{1},\cdots,\tilde{v}_{s},\tilde{v}_{1}^{\ast}%
,\cdots,\tilde{v}_{s}^{\ast}\right\rangle  &  =\left\vert v_{1},\cdots
,v_{s},v_{1}^{\ast},\cdots,v_{s}^{\ast}\right\rangle \left(
\begin{array}
[c]{cc}%
-\boldsymbol{\alpha}^{\frac{1}{2}} & \mathbb{O}\\
\mathbb{O} & \boldsymbol{\alpha}^{\frac{1}{2}}%
\end{array}
\right)  .
\end{align}
The real invariance group associated with $A$ is given by
\begin{align}
TAT^{\dagger}  &  =A\\
R_{s}\left(  T\right)   &  \in\mathcal{B}\left(  \mathbb{R}^{2s}\right)
\end{align}
giving
\begin{equation}
R_{B}\left(  T\right)  R_{B}\left(  A\right)  R_{B}\left(  T\right)
^{\dagger}=R_{B}\left(  A\right)  ,
\end{equation}
where $R_{B}\left(  T\right)  $ and $R_{B}\left(  A\right)  $ are the
representations wrt the basis $B.$ This leads to
\begin{equation}
R_{s}\left(  T\right)  \left(
\begin{array}
[c]{cc}%
\mathbb{O} & \boldsymbol{\alpha}\\
-\boldsymbol{\alpha} & \mathbb{O}%
\end{array}
\right)  R_{s}\left(  T\right)  ^{t}=\left(
\begin{array}
[c]{cc}%
\mathbb{O} & \boldsymbol{\alpha}\\
-\boldsymbol{\alpha} & \mathbb{O}%
\end{array}
\right)  ,
\end{equation}%
\begin{equation}
R_{v}\left(  T\right)  \left(
\begin{array}
[c]{cc}%
-\boldsymbol{\alpha} & \mathbb{O}\\
\mathbb{O} & \boldsymbol{\alpha}%
\end{array}
\right)  R_{v}\left(  T\right)  ^{\dagger}=\left(
\begin{array}
[c]{cc}%
-\boldsymbol{\alpha} & \mathbb{O}\\
\mathbb{O} & \boldsymbol{\alpha}%
\end{array}
\right)
\end{equation}
and
\begin{align}
T  &  =\left\vert e_{1},\cdots\cdots,e_{2s}\right\rangle R_{e}\left(
T\right)  \left\langle e_{1},\cdots\cdots,e_{2s}\right\vert =\left\vert
e_{1},\cdots\cdots,e_{2s}\right\rangle \mathbb{O}R_{e}\left(  T\right)
\mathbb{O}^{t}\left\langle e_{1},\cdots\cdots,e_{2s}\right\vert \\
&  =\left\vert s_{1},\cdots\cdots,s_{2s}\right\rangle R_{s}\left(  T\right)
\left\langle s_{1},\cdots\cdots,s_{2s}\right\vert =\left\vert e_{1}%
,\cdots\cdots,e_{2s}\right\rangle \mathbb{V}R_{v}\left(  T\right)
\mathbb{V}^{\dagger}\left\langle e_{1},\cdots\cdots,e_{2s}\right\vert \\
&  =\left\vert v_{1},\cdots,v_{s},v_{1}^{\ast},\cdots,v_{s}^{\ast
}\right\rangle R_{v}\left(  T\right)  \left\langle v_{1},\cdots,v_{s}%
,v_{1}^{\ast},\cdots,v_{s}^{\ast}\right\vert \\
&  =\left\vert s_{1},\cdots\cdots,s_{2s}\right\rangle \frac{1}{\sqrt{2}%
}\left(
\begin{array}
[c]{cc}%
\mathbb{I}_{s} & \mathbb{I}_{s}\\
i\mathbb{I}_{s} & -i\mathbb{I}_{s}%
\end{array}
\right)  R_{v}\left(  T\right)  \frac{1}{\sqrt{2}}\left(
\begin{array}
[c]{cc}%
\mathbb{I}_{s} & -i\mathbb{I}_{s}\\
\mathbb{I}_{s} & i\mathbb{I}_{s}%
\end{array}
\right)  \left\langle s_{1},\cdots\cdots,s_{2s}\right\vert .
\end{align}


\subsubsection{Real subspace of complex space}

Note that $\mathbb{R}^{2s}$ can be embedded in $\mathbb{C}^{2s}$ as
\begin{equation}
\mathbb{R}^{2s}\hookrightarrow\mathbb{C}^{2s}=\mathbb{C}^{s}\oplus
\overline{\mathbb{C}^{s}},
\end{equation}
while a complexification of $\mathbb{R}^{2s}$ produces $\mathbb{C}^{s}.$ This
can be abstractly expressed as
\begin{equation}
V_{2s\mathbb{R}}\hookrightarrow V_{s\mathbb{C}}\oplus\overline{V_{s\mathbb{C}%
}}%
\end{equation}
where a complexification of $V_{2s\mathbb{R}}$ produces $V_{s\mathbb{C}}$.

As%
\begin{equation}
\left\vert v_{1},\cdots,v_{s},v_{1}^{\ast},\cdots,v_{s}^{\ast}\right\rangle
=\left\vert s_{1},\cdots\cdots,s_{2s}\right\rangle \frac{1}{\sqrt{2}}\left(
\begin{array}
[c]{cc}%
\mathbb{I}_{s} & \mathbb{I}_{s}\\
i\mathbb{I}_{s} & -i\mathbb{I}_{s}%
\end{array}
\right)
\end{equation}
one obtains%
\begin{align}
R_{s}\left(  T\right)   &  =\frac{1}{\sqrt{2}}\left(
\begin{array}
[c]{cc}%
\mathbb{I}_{s} & \mathbb{I}_{s}\\
i\mathbb{I}_{s} & -i\mathbb{I}_{s}%
\end{array}
\right)  R_{v}\left(  T\right)  \frac{1}{\sqrt{2}}\left(
\begin{array}
[c]{cc}%
\mathbb{I}_{s} & -i\mathbb{I}_{s}\\
\mathbb{I}_{s} & i\mathbb{I}_{s}%
\end{array}
\right) \\
R_{v}\left(  T\right)   &  =\frac{1}{\sqrt{2}}\left(
\begin{array}
[c]{cc}%
\mathbb{I}_{s} & -i\mathbb{I}_{s}\\
\mathbb{I}_{s} & i\mathbb{I}_{s}%
\end{array}
\right)  R_{s}\left(  T\right)  \frac{1}{\sqrt{2}}\left(
\begin{array}
[c]{cc}%
\mathbb{I}_{s} & \mathbb{I}_{s}\\
i\mathbb{I}_{s} & -i\mathbb{I}_{s}%
\end{array}
\right)  .
\end{align}
$\lceil$

The real representation in terms of the complex one is
\begin{align}
R_{s}\left(  T\right)   &  =\left(
\begin{array}
[c]{cc}%
T_{s11} & T_{s12}\\
T_{s21} & T_{s22}%
\end{array}
\right)  =\frac{1}{\sqrt{2}}\left(
\begin{array}
[c]{cc}%
\mathbb{I}_{s} & \mathbb{I}_{s}\\
i\mathbb{I}_{s} & -i\mathbb{I}_{s}%
\end{array}
\right)  \left(
\begin{array}
[c]{cc}%
T_{v11} & T_{v12}\\
T_{v21} & T_{v22}%
\end{array}
\right)  \frac{1}{\sqrt{2}}\left(
\begin{array}
[c]{cc}%
\mathbb{I}_{s} & -i\mathbb{I}_{s}\\
\mathbb{I}_{s} & i\mathbb{I}_{s}%
\end{array}
\right) \\
&  =\frac{1}{2}\left(
\begin{array}
[c]{cc}%
T_{v11}+T_{v21} & T_{v12}+T_{v22}\\
i\left(  T_{v11}-T_{v21}\right)  & i\left(  T_{v12}-T_{v22}\right)
\end{array}
\right)  \left(
\begin{array}
[c]{cc}%
\mathbb{I}_{s} & -i\mathbb{I}_{s}\\
\mathbb{I}_{s} & i\mathbb{I}_{s}%
\end{array}
\right) \\
&  =\frac{1}{2}\left(
\begin{array}
[c]{cc}%
T_{v11}+T_{v21}+T_{v12}+T_{v22} & -i\left(  T_{v11}+T_{v21}-T_{v12}%
-T_{v22}\right) \\
i\left(  T_{v11}-T_{v21}+T_{v12}-T_{v22}\right)  & T_{v11}-T_{v21}%
-T_{v12}+T_{v22}%
\end{array}
\right) \\
&  =\frac{1}{2}\left(
\begin{array}
[c]{cc}%
T_{v11}+T_{v22}+T_{v21}+T_{v12} & -i\left(  T_{v11}-T_{v22}+T_{v21}%
-T_{v12}\right) \\
i\left(  T_{v11}-T_{v22}-T_{v21}+T_{v12}\right)  & T_{v11}+T_{v22}%
-T_{v21}-T_{v12}%
\end{array}
\right)  ,
\end{align}
while the complex one in terms of the real one is
\begin{align*}
R_{v}\left(  T\right)   &  =\left(
\begin{array}
[c]{cc}%
T_{v11} & T_{v12}\\
T_{v21} & T_{v22}%
\end{array}
\right)  =\frac{1}{\sqrt{2}}\left(
\begin{array}
[c]{cc}%
\mathbb{I}_{s} & -i\mathbb{I}_{s}\\
\mathbb{I}_{s} & i\mathbb{I}_{s}%
\end{array}
\right)  \left(
\begin{array}
[c]{cc}%
T_{s11} & T_{s12}\\
T_{s21} & T_{s22}%
\end{array}
\right)  \frac{1}{\sqrt{2}}\left(
\begin{array}
[c]{cc}%
\mathbb{I}_{s} & \mathbb{I}_{s}\\
i\mathbb{I}_{s} & -i\mathbb{I}_{s}%
\end{array}
\right) \\
&  =\frac{1}{2}\left(
\begin{array}
[c]{cc}%
T_{s11}-iT_{s21} & T_{s12}-iT_{s22}\\
T_{s11}+iT_{s21} & T_{s12}+iT_{s22}%
\end{array}
\right)  \left(
\begin{array}
[c]{cc}%
\mathbb{I}_{s} & \mathbb{I}_{s}\\
i\mathbb{I}_{s} & -i\mathbb{I}_{s}%
\end{array}
\right) \\
&  =\frac{1}{2}\left(
\begin{array}
[c]{cc}%
T_{s11}-iT_{s21}+i\left(  T_{s12}-iT_{s22}\right)  & T_{s11}-iT_{s21}-i\left(
T_{s12}-iT_{s22}\right) \\
T_{s11}+iT_{s21}+i\left(  T_{s12}+iT_{s22}\right)  & T_{s11}+iT_{s21}-i\left(
T_{s12}+iT_{s22}\right)
\end{array}
\right) \\
&  =\frac{1}{2}\left(
\begin{array}
[c]{cc}%
T_{s11}-iT_{s21}+iT_{s12}+T_{s22} & T_{s11}-iT_{s21}-iT_{s12}-T_{s22}\\
T_{s11}+iT_{s21}+iT_{s12}-T_{s22} & T_{s11}+iT_{s21}-iT_{s12}+T_{s22}%
\end{array}
\right) \\
&  =\frac{1}{2}\left(
\begin{array}
[c]{cc}%
T_{s11}+T_{s22}-iT_{s21}+iT_{s12} & T_{s11}-T_{s22}-iT_{s21}-iT_{s12}\\
T_{s11}-T_{s22}+iT_{s21}+iT_{s12} & T_{s11}+T_{s22}+iT_{s21}-iT_{s12}%
\end{array}
\right) \\
&  =\frac{1}{2}\left(
\begin{array}
[c]{cc}%
T_{s11}+T_{s22}-iT_{s21}+iT_{s12} & T_{s11}-T_{s22}-iT_{s21}-iT_{s12}\\
\left(  T_{s11}-T_{s22}-iT_{s21}-iT_{s12}\right)  ^{\ast} & \left(
T_{s11}+T_{s22}-iT_{s21}+iT_{s12}\right)  ^{\ast}%
\end{array}
\right)
\end{align*}
i.e.
\begin{equation}
R_{v}\left(  T\right)  =\left(
\begin{array}
[c]{cc}%
T_{v11} & T_{v12}\\
\left(  T_{v12}\right)  ^{\ast} & \left(  T_{v11}\right)  ^{\ast}%
\end{array}
\right)  .
\end{equation}
Using this in the preceding expression for $R_{s}\left(  T\right)  $ gives
\begin{align}
R_{s}\left(  T\right)   &  =\frac{1}{2}\left(
\begin{array}
[c]{cc}%
T_{v11}+\left(  T_{v11}\right)  ^{\ast}+\left(  T_{v12}\right)  ^{\ast
}+T_{v12} & -i\left(  T_{v11}-\left(  T_{v11}\right)  ^{\ast}-\left(
T_{v12}-\left(  T_{v12}\right)  ^{\ast}\right)  \right) \\
i\left(  T_{v11}-\left(  T_{v11}\right)  ^{\ast}+T_{v12}-\left(
T_{v12}\right)  ^{\ast}\right)  & T_{v11}+\left(  T_{v11}\right)  ^{\ast
}-\left(  T_{v12}+\left(  T_{v12}\right)  ^{\ast}\right)
\end{array}
\right) \\
&  =\left(
\begin{array}
[c]{cc}%
\operatorname{Re}\left(  T_{v11}+T_{v12}\right)  & \operatorname{Im}\left(
T_{v11}-T_{v12}\right) \\
-\operatorname{Im}\left(  T_{v11}+T_{v12}\right)  & \operatorname{Re}\left(
T_{v11}-T_{v12}\right)
\end{array}
\right)  .
\end{align}
$\lceil$ This can also be seen from%
\begin{align}
R_{s}\left(  T\right)   &  =\frac{1}{\sqrt{2}}\left(
\begin{array}
[c]{cc}%
\mathbb{I}_{s} & \mathbb{I}_{s}\\
i\mathbb{I}_{s} & -i\mathbb{I}_{s}%
\end{array}
\right)  \left(
\begin{array}
[c]{cc}%
T_{v11} & T_{v12}\\
\overline{T_{v12}} & \overline{T_{v11}}%
\end{array}
\right)  \frac{1}{\sqrt{2}}\left(
\begin{array}
[c]{cc}%
\mathbb{I}_{s} & -i\mathbb{I}_{s}\\
\mathbb{I}_{s} & i\mathbb{I}_{s}%
\end{array}
\right) \\
&  =\frac{1}{2}\left(
\begin{array}
[c]{cc}%
T_{v11}+\overline{T_{v12}} & T_{v12}+\overline{T_{v11}}\\
iT_{v11}-i\overline{T_{v12}} & iT_{v12}-i\overline{T_{v11}}%
\end{array}
\right)  \left(
\begin{array}
[c]{cc}%
\mathbb{I}_{s} & -i\mathbb{I}_{s}\\
\mathbb{I}_{s} & i\mathbb{I}_{s}%
\end{array}
\right) \\
&  =\frac{1}{2}\left(
\begin{array}
[c]{cc}%
T_{v11}+\overline{T_{v12}}+T_{v12}+\overline{T_{v11}} & -i\left(
T_{v11}+\overline{T_{v12}}\right)  +i\left(  T_{v12}+\overline{T_{v11}}\right)
\\
i\left(  T_{v11}-\overline{T_{v12}}\right)  +i\left(  T_{v12}-\overline
{T_{v11}}\right)  & \left(  T_{v11}-\overline{T_{v12}}\right)  -\left(
T_{v12}-\overline{T_{v11}}\right)
\end{array}
\right) \\
&  =\frac{1}{2}\left(
\begin{array}
[c]{cc}%
2\operatorname{Re}\left(  T_{v11}+T_{v12}\right)  & +i\left(  T_{v12}%
-\overline{T_{v12}}\right)  +i\left(  \overline{T_{v11}}-T_{v11}\right) \\
i\left(  T_{v11}-\overline{T_{v11}}\right)  +i\left(  T_{v12}-\overline
{T_{v12}}\right)  & 2\operatorname{Re}\left(  T_{v11}-T_{v12}\right)
\end{array}
\right) \\
&  =\left(
\begin{array}
[c]{cc}%
\operatorname{Re}\left(  T_{v11}+T_{v12}\right)  & \operatorname{Im}\left(
T_{v11}-T_{v12}\right) \\
-\operatorname{Im}\left(  T_{v11}+T_{v12}\right)  & \operatorname{Re}\left(
T_{v11}-T_{v12}\right)
\end{array}
\right)
\end{align}
$\rfloor$

$\lceil$Note that if $r=2$
\begin{equation}
R_{s}\left(  T\right)  \left(
\begin{array}
[c]{cc}%
0 & 1\\
-1 & 0
\end{array}
\right)  R_{s}\left(  T\right)  ^{t}=\left(
\begin{array}
[c]{cc}%
0 & 1\\
-1 & 0
\end{array}
\right)
\end{equation}
and
\begin{equation}
R_{v}\left(  T\right)  \left(
\begin{array}
[c]{cc}%
-1 & 0\\
0 & 1
\end{array}
\right)  R_{v}\left(  T\right)  ^{\dagger}=\left(
\begin{array}
[c]{cc}%
-1 & 0\\
0 & 1
\end{array}
\right)
\end{equation}
i.e. $T\in Sp\left(  2,\mathbb{R}\right)  $ and $T\in U\left(  1,1\right)
.\rfloor$

$\lceil$If $r=4$ this gives
\begin{equation}
R_{s}\left(  T\right)  \left(
\begin{array}
[c]{cccc}%
0 & 0 & \alpha_{1} & 0\\
0 & 0 & 0 & \alpha_{2}\\
-\alpha_{1} & 0 & 0 & 0\\
0 & -\alpha_{2} & 0 & 0
\end{array}
\right)  R_{s}\left(  T\right)  ^{t}=\left(
\begin{array}
[c]{cccc}%
0 & 0 & \alpha_{1} & 0\\
0 & 0 & 0 & \alpha_{2}\\
-\alpha_{1} & 0 & 0 & 0\\
0 & -\alpha_{2} & 0 & 0
\end{array}
\right)
\end{equation}
and
\begin{equation}
R_{v}\left(  T\right)  \left(
\begin{array}
[c]{cccc}%
-\alpha_{1} & 0 & 0 & 0\\
0 & -\alpha_{2} & 0 & 0\\
0 & 0 & \alpha_{1} & 0\\
0 & 0 & 0 & \alpha_{2}%
\end{array}
\right)  R_{v}\left(  T\right)  ^{\dagger}=\left(
\begin{array}
[c]{cccc}%
-\alpha_{1} & 0 & 0 & 0\\
0 & -\alpha_{2} & 0 & 0\\
0 & 0 & \alpha_{1} & 0\\
0 & 0 & 0 & \alpha_{2}%
\end{array}
\right)  .
\end{equation}
If $\alpha_{1}\neq\alpha_{2}$ then $T\in Sp\left(  2,\mathbb{R}\right)  \times
Sp\left(  2,\mathbb{R}\right)  $ and $T\in U\left(  1,1\right)  \times
U\left(  1,1\right)  ,$ while if $\alpha_{1}=\alpha_{2}$ then $T\in Sp\left(
4,\mathbb{R}\right)  $ and $T\in U\left(  2,2\right)  .\rfloor$

If $R_{s}\left(  T\right)  $ corresponds to a complex transformation
\begin{equation}
R_{\mathbb{C}s}\left(  T\right)  :\mathbb{C}^{s}\rightarrow\mathbb{C}^{s}%
\end{equation}
then
\begin{equation}
\left(
\begin{array}
[c]{cc}%
t_{R} & -t_{I}\\
t_{I} & t_{R}%
\end{array}
\right)  =\left(
\begin{array}
[c]{cc}%
\operatorname{Re}\left(  T_{v11}+T_{v12}\right)  & \operatorname{Im}\left(
T_{v11}-T_{v12}\right) \\
-\operatorname{Im}\left(  T_{v11}+T_{v12}\right)  & \operatorname{Re}\left(
T_{v11}-T_{v12}\right)
\end{array}
\right)  ,
\end{equation}
which implies that
\begin{equation}
T_{v12}=\mathbb{O}%
\end{equation}
i.e.
\begin{align}
R_{s}\left(  T\right)   &  =\left(
\begin{array}
[c]{cc}%
\operatorname{Re}T_{v11} & \operatorname{Im}T_{v11}\\
-\operatorname{Im}T_{v11} & \operatorname{Re}T_{v11}%
\end{array}
\right) \\
R_{v}\left(  T\right)   &  =\left(
\begin{array}
[c]{cc}%
T_{v11} & \mathbb{O}\\
\mathbb{O} & \overline{T_{v11}}%
\end{array}
\right)  .
\end{align}


\paragraph{Orthogonal}%

\begin{align}
\mathbb{B}  &  =\mathbb{O}\left(
\begin{array}
[c]{cc}%
\boldsymbol{\alpha} & \mathbb{O}\\
\mathbb{O} & \boldsymbol{\alpha}%
\end{array}
\right)  \mathbb{O}^{t}=\left\vert \mathbf{s}_{1},\cdots,\mathbf{s}%
_{\mathbf{s}},\mathbf{s}_{\mathbf{s}+1},\cdots,\mathbf{s}_{\mathbf{s}%
+\mathbf{s}}\right\rangle \left(
\begin{array}
[c]{cc}%
\boldsymbol{\alpha} & \mathbb{O}\\
\mathbb{O} & \boldsymbol{\alpha}%
\end{array}
\right)  \left\langle \mathbf{s}_{1},\cdots,\mathbf{s}_{\mathbf{s}}%
,\mathbf{s}_{\mathbf{s}+1},\cdots,\mathbf{s}_{\mathbf{s}+\mathbf{s}%
}\right\vert \\
&  =\frac{1}{2}\left\vert \mathbf{v}_{1},\cdots,\mathbf{v}_{\mathbf{s}%
},\mathbf{v}_{1}^{\ast},\cdots,\mathbf{v}_{\mathbf{s}}^{\ast}\right\rangle
\left(
\begin{array}
[c]{cc}%
\mathbb{I}_{s} & -i\mathbb{I}_{s}\\
\mathbb{I}_{s} & i\mathbb{I}_{s}%
\end{array}
\right)  \left(
\begin{array}
[c]{cc}%
\boldsymbol{\alpha} & \mathbb{O}\\
\mathbb{O} & \boldsymbol{\alpha}%
\end{array}
\right)  \left(
\begin{array}
[c]{cc}%
\mathbb{I}_{s} & \mathbb{I}_{s}\\
i\mathbb{I}_{s} & -i\mathbb{I}_{s}%
\end{array}
\right)  \left\langle \mathbf{v}_{1},\cdots,\mathbf{v}_{\mathbf{s}}%
,\mathbf{v}_{1}^{\ast},\cdots,\mathbf{v}_{\mathbf{s}}^{\ast}\right\vert \\
&  =\left\vert \mathbf{v}_{1},\cdots,\mathbf{v}_{\mathbf{s}},\mathbf{v}%
_{1}^{\ast},\cdots,\mathbf{v}_{\mathbf{s}}^{\ast}\right\rangle \left(
\begin{array}
[c]{cc}%
\boldsymbol{\alpha} & \mathbb{O}\\
\mathbb{O} & \boldsymbol{\alpha}%
\end{array}
\right)  \left\langle \mathbf{v}_{1},\cdots,\mathbf{v}_{\mathbf{s}}%
,\mathbf{v}_{1}^{\ast},\cdots,\mathbf{v}_{\mathbf{s}}^{\ast}\right\vert
\end{align}
i.e.%
\begin{equation}
\frac{1}{2}\left(
\begin{array}
[c]{cc}%
\mathbb{I}_{s} & -i\mathbb{I}_{s}\\
\mathbb{I}_{s} & i\mathbb{I}_{s}%
\end{array}
\right)  \left(
\begin{array}
[c]{cc}%
\boldsymbol{\alpha} & \mathbb{O}\\
\mathbb{O} & \boldsymbol{\alpha}%
\end{array}
\right)  \left(
\begin{array}
[c]{cc}%
\mathbb{I}_{s} & \mathbb{I}_{s}\\
i\mathbb{I}_{s} & -i\mathbb{I}_{s}%
\end{array}
\right)  =\left(
\begin{array}
[c]{cc}%
\boldsymbol{\alpha} & \mathbb{O}\\
\mathbb{O} & \boldsymbol{\alpha}%
\end{array}
\right)  .
\end{equation}
$\lceil$%
\begin{align}
\frac{1}{2}\left(
\begin{array}
[c]{cc}%
\mathbb{I}_{s} & -i\mathbb{I}_{s}\\
\mathbb{I}_{s} & i\mathbb{I}_{s}%
\end{array}
\right)  \left(
\begin{array}
[c]{cc}%
\boldsymbol{\alpha} & \mathbb{O}\\
\mathbb{O} & \boldsymbol{\alpha}%
\end{array}
\right)  \left(
\begin{array}
[c]{cc}%
\mathbb{I}_{s} & \mathbb{I}_{s}\\
i\mathbb{I}_{s} & -i\mathbb{I}_{s}%
\end{array}
\right)   &  =\frac{1}{2}\left(
\begin{array}
[c]{cc}%
\boldsymbol{\alpha} & -i\boldsymbol{\alpha}\\
\boldsymbol{\alpha} & i\boldsymbol{\alpha}%
\end{array}
\right)  \left(
\begin{array}
[c]{cc}%
\mathbb{I}_{s} & \mathbb{I}_{s}\\
i\mathbb{I}_{s} & -i\mathbb{I}_{s}%
\end{array}
\right) \\
&  =\frac{1}{2}\left(
\begin{array}
[c]{cc}%
2\boldsymbol{\alpha} & \mathbb{O}\\
\mathbb{O} & 2\boldsymbol{\alpha}%
\end{array}
\right)  =\left(
\begin{array}
[c]{cc}%
\boldsymbol{\alpha} & \mathbb{O}\\
\mathbb{O} & \boldsymbol{\alpha}%
\end{array}
\right)  .
\end{align}


$\rfloor$

Thus
\begin{equation}
B=\left\vert e_{1},\cdots\cdots,e_{2\mathbf{s}}\right\rangle \mathbb{O}\left(
%
\begin{array}
[c]{cc}%
\boldsymbol{\alpha} & \mathbb{O}\\
\mathbb{O} & \boldsymbol{\alpha}%
\end{array}
\right)  \mathbb{O}^{t}\left\langle e_{1},\cdots\cdots,e_{2\mathbf{s}%
}\right\vert =\left\vert e_{1},\cdots\cdots,e_{2\mathbf{s}}\right\rangle
\mathbb{V}\left(
\begin{array}
[c]{cc}%
\boldsymbol{\alpha} & \mathbb{O}\\
\mathbb{O} & \boldsymbol{\alpha}%
\end{array}
\right)  \mathbb{V}^{\dagger}\left\langle e_{1},\cdots\cdots,e_{2\mathbf{s}%
}\right\vert .
\end{equation}
The invariance group associated with $B$ is given by
\begin{equation}
TBT^{\dagger}=B
\end{equation}
giving
\begin{equation}
R_{B}\left(  T\right)  R_{B}\left(  B\right)  R_{B}\left(  T\right)
^{\dagger}=R_{B}\left(  B\right)  .
\end{equation}


\subsection{Invariances}

The real antisymmetric form is defined by
\begin{equation}
A=-i\sum_{1\leq j\leq s}\alpha_{j}\left\{  \left\vert v_{j}\right\rangle
\left\langle v_{j}\right\vert -\left\vert v_{j}^{\ast}\right\rangle
\left\langle v_{j}^{\ast}\right\vert \right\}  ;\qquad\alpha_{j}\in
\mathbb{R}^{+},
\end{equation}
which corresponds to the geminal%
\begin{equation}
\left\vert g\right\rangle =-i\sum_{1\leq j\leq s}\alpha_{j}\left\vert
v_{j}v_{j}^{\ast}\right\rangle .
\end{equation}
One can define the orthogonal projectors
\begin{align}
P_{A+}  &  =\sum_{1\leq j\leq s}\left\vert v_{j}\right\rangle \left\langle
v_{j}\right\vert \\
P_{A-}  &  =\sum_{1\leq j\leq s}\left\vert v_{j}^{\ast}\right\rangle
\left\langle v_{j}^{\ast}\right\vert
\end{align}
so that
\begin{equation}
\left[  A,P_{A+}\right]  =\left[  A,P_{A+}\right]  =0
\end{equation}
and
\begin{equation}
A=P_{A+}AP_{A+}+P_{A-}AP_{A-}.
\end{equation}


If
\begin{equation}
\left\vert w_{j},w_{j}^{\ast}\right\rangle =\left\vert v_{j},v_{j}^{\ast
}\right\rangle \left(
\begin{array}
[c]{cc}%
\alpha & b\\
b^{\ast} & a^{\ast}%
\end{array}
\right)  =\left\vert v_{j},v_{j}^{\ast}\right\rangle T,
\end{equation}
where
\begin{equation}
\det\left(
\begin{array}
[c]{cc}%
a & b\\
b^{\ast} & a^{\ast}%
\end{array}
\right)  =1=\left\vert a\right\vert ^{2}-\left\vert b\right\vert ^{2}%
\end{equation}
and
\begin{equation}
\left(
\begin{array}
[c]{cc}%
\alpha & b\\
b^{\ast} & a^{\ast}%
\end{array}
\right)  \left(
\begin{array}
[c]{cc}%
\alpha & b\\
b^{\ast} & a^{\ast}%
\end{array}
\right)  ^{\dagger}=\left(
\begin{array}
[c]{cc}%
\alpha & b\\
b^{\ast} & a^{\ast}%
\end{array}
\right)  \left(
\begin{array}
[c]{cc}%
\alpha^{\ast} & b\\
b^{\ast} & a
\end{array}
\right)  =\left(
\begin{array}
[c]{cc}%
\left\vert a\right\vert ^{2}+\left\vert b\right\vert ^{2} & ab+\left(
ab\right)  ^{\ast}\\
ab+\left(  ab\right)  ^{\ast} & \left\vert a\right\vert ^{2}+\left\vert
b\right\vert ^{2}%
\end{array}
\right)  =\left(
\begin{array}
[c]{cc}%
1 & 0\\
0 & 1
\end{array}
\right)
\end{equation}
then%
\begin{align}
\left\vert a\right\vert ^{2}+\left\vert b\right\vert ^{2}  &  =1\\
\operatorname{Re}\left(  ab\right)   &  =0
\end{align}%
\begin{gather}
1=\left\vert a\right\vert ^{2}-\left\vert b\right\vert ^{2}\\
1=\left\vert a\right\vert ^{2}+\left\vert b\right\vert ^{2}\\
\Downarrow\\
a=e^{i\theta};\quad b=0
\end{gather}
i.e.%
\begin{equation}
T=\left(
\begin{array}
[c]{cc}%
e^{i\theta} & 0\\
0 & e^{-i\theta}%
\end{array}
\right)  \Rightarrow T\in SU\left(  2\right)  \cap SU\left(  1,1\right)  .
\end{equation}


\begin{proposition}
If $\alpha_{1}=\alpha_{2}=\cdots=\alpha_{d}$ and
\begin{equation}
\left\vert w_{1},\cdots,w_{d},w_{1}^{\ast},\cdots,w_{d}^{\ast}\right\rangle
=T\left\vert v_{1},\cdots,v_{d},v_{1}^{\ast},\cdots,v_{d}^{\ast}\right\rangle
=\left\vert v_{1},\cdots,v_{d},v_{1}^{\ast},\cdots,v_{d}^{\ast}\right\rangle
\left(
\begin{array}
[c]{cc}%
\mathbb{A} & \mathbb{B}\\
\mathbb{B}^{\ast} & \mathbb{A}^{\ast}%
\end{array}
\right)  =\left\vert v_{j},\cdots,v_{j}^{\ast}\right\rangle \mathbb{T}%
\end{equation}
and $T\in SU\left(  2d\right)  $ \textbf{then} $T\in SU\left(  d,d\right)  $
i.e. $T\in SU\left(  2d\right)  \cap SU\left(  d,d\right)  .$
\end{proposition}

\begin{proof}%
\begin{align}
\left(
\begin{array}
[c]{cc}%
\mathbb{I}_{d} & \mathbb{O}\\
\mathbb{O} & \mathbb{I}_{d}%
\end{array}
\right)   &  =\left(
\begin{array}
[c]{cc}%
\mathbb{A} & \mathbb{B}\\
\mathbb{B}^{\ast} & \mathbb{A}^{\ast}%
\end{array}
\right)  \left(
\begin{array}
[c]{cc}%
\mathbb{A} & \mathbb{B}\\
\mathbb{B}^{\ast} & \mathbb{A}^{\ast}%
\end{array}
\right)  ^{\dagger}=\left(
\begin{array}
[c]{cc}%
\mathbb{A} & \mathbb{B}\\
\mathbb{B}^{\ast} & \mathbb{A}^{\ast}%
\end{array}
\right)  \left(
\begin{array}
[c]{cc}%
\mathbb{A}^{\dagger} & \mathbb{B}^{t}\\
\mathbb{B}^{\dagger} & \mathbb{A}^{t}%
\end{array}
\right) \\
&  =\left(
\begin{array}
[c]{cc}%
\mathbb{AA}^{\dagger}+\mathbb{BB}^{\dagger} & \mathbb{AB}^{t}+\mathbb{BA}%
^{t}\\
\mathbb{B}^{\ast}\mathbb{A}^{\dagger}+\mathbb{A}^{\ast}\mathbb{B}^{\dagger} &
\mathbb{A}^{\ast}\mathbb{A}^{t}+\mathbb{B}^{\ast}\mathbb{B}^{t}%
\end{array}
\right)
\end{align}
together with
\begin{align}
\left(
\begin{array}
[c]{cc}%
\mathbb{I}_{d} & \mathbb{O}\\
\mathbb{O} & -\mathbb{I}_{d}%
\end{array}
\right)   &  =\left(
\begin{array}
[c]{cc}%
\mathbb{A} & \mathbb{B}\\
\mathbb{B}^{\ast} & \mathbb{A}^{\ast}%
\end{array}
\right)  \left(
\begin{array}
[c]{cc}%
\mathbb{I}_{d} & \mathbb{O}\\
\mathbb{O} & -\mathbb{I}_{d}%
\end{array}
\right)  \left(
\begin{array}
[c]{cc}%
\mathbb{A} & \mathbb{B}\\
\mathbb{B}^{\ast} & \mathbb{A}^{\ast}%
\end{array}
\right)  ^{\dagger}=\left(
\begin{array}
[c]{cc}%
\mathbb{A} & -\mathbb{B}\\
\mathbb{B}^{\ast} & -\mathbb{A}^{\ast}%
\end{array}
\right)  \left(
\begin{array}
[c]{cc}%
\mathbb{A}^{\dagger} & \mathbb{B}^{t}\\
\mathbb{B}^{\dagger} & \mathbb{A}^{t}%
\end{array}
\right) \\
&  =\left(
\begin{array}
[c]{cc}%
\mathbb{AA}^{\dagger}-\mathbb{BB}^{\dagger} & \mathbb{AB}^{t}-\mathbb{BA}%
^{t}\\
\mathbb{B}^{\ast}\mathbb{A}^{\dagger}-\mathbb{A}^{\ast}\mathbb{B}^{\dagger} &
-\mathbb{A}^{\ast}\mathbb{A}^{t}+\mathbb{B}^{\ast}\mathbb{B}^{t}%
\end{array}
\right)
\end{align}
implies that
\begin{align}
\mathbb{T}  &  \mathbb{=}\left(
\begin{array}
[c]{cc}%
\mathbb{A} & \mathbb{O}\\
\mathbb{O} & \mathbb{A}^{\ast}%
\end{array}
\right) \\
\mathbb{AA}^{\dagger}  &  =\mathbb{I}_{d}.
\end{align}
Hence
\begin{equation}
A=-i\sum_{1\leq k\leq\kappa}\alpha_{\kappa}\sum_{1\leq j\leq d_{j}}\left\{
\left\vert v_{\kappa j}\right\rangle \left\langle v_{\kappa j}\right\vert
-\left\vert \left(  v_{\kappa j}\right)  ^{\ast}\right\rangle \left\langle
\left(  v_{\kappa j}\right)  ^{\ast}\right\vert \right\}
\end{equation}
giving%
\begin{equation}
\left\vert g\right\rangle =-i\sum_{1\leq k\leq\kappa}\alpha_{\kappa}%
\sum_{1\leq j\leq d_{j}}\left\vert \left(  v_{\kappa j}\right)  \left(
v_{\kappa j}\right)  ^{\ast}\right\rangle .
\end{equation}

\end{proof}

\subsection{Antisymmetric Structure}

Antisymmetric bilinear forms
\begin{equation}
g:V\times V\rightarrow\mathbb{F}%
\end{equation}
can be expressed as
\begin{equation}
g\left(  a,b\right)  =\left\langle g\left\vert a\wedge b\right.  \right\rangle
=\left\langle a\left\vert Ab\right.  \right\rangle =\mathbf{a}^{t}%
\mathbb{A}\mathbf{b},
\end{equation}
where
\begin{align}
\left\langle g\right\vert  &  \in\left(  V\wedge V\right)  ^{d}\\
\left\langle a\right\vert  &  \in V^{d}\\
A &  \in\mathcal{B}\left(  V\right)  \\
a,b &  \in V
\end{align}
$\lceil$There is an injective map
\begin{equation}
\Upsilon:V\rightarrow V^{d}%
\end{equation}
by considering biorthogonal bases$\rfloor$

The cannonical form of a real geminal is
\begin{equation}
\left\vert g\right\rangle =\sum\limits_{1\leq j\leq s}\alpha_{j}\left\vert
s_{j}s_{j+s}\right\rangle =-i\sum_{1\leq j\leq s}\alpha_{j}\left\vert
v_{j}v_{j}^{\ast}\right\rangle .
\end{equation}


\bigskip%

\begin{equation}
A=\left\vert s_{1},\cdots,s_{s},s_{s+1},\cdots,s_{s+s}\right\rangle \left(
\begin{array}
[c]{cc}%
\mathbb{O} & \boldsymbol{\alpha}\\
-\boldsymbol{\alpha} & \mathbb{O}%
\end{array}
\right)  \left\langle s_{1},\cdots,s_{s},s_{s+1},\cdots,s_{s+s}\right\vert
=\left\vert v_{1},\cdots,v_{s},v_{1}^{\ast},\cdots,v_{s}^{\ast}\right\rangle
i\left(
\begin{array}
[c]{cc}%
-\boldsymbol{\alpha} & \mathbb{O}\\
\mathbb{O} & \boldsymbol{\alpha}%
\end{array}
\right)  \left\langle v_{1},\cdots,v_{\mathbf{s}},v_{1}^{\ast},\cdots
,v_{s}^{\ast}\right\vert
\end{equation}%
\begin{equation}
J_{A}=\left\vert s_{1},\cdots,s_{s},s_{s+1},\cdots,s_{s+s}\right\rangle
\left(
\begin{array}
[c]{cc}%
\mathbb{O} & \mathbb{I}_{s}\\
-\mathbb{I}_{s} & \mathbb{O}%
\end{array}
\right)  \left\langle s_{1},\cdots,s_{s},s_{s+1},\cdots,s_{s+s}\right\vert
=\left\vert v_{1},\cdots,v_{s},v_{1}^{\ast},\cdots,v_{s}^{\ast}\right\rangle
i\left(
\begin{array}
[c]{cc}%
-\mathbb{I}_{s} & \mathbb{O}\\
\mathbb{O} & \mathbb{I}_{s}%
\end{array}
\right)  \left\langle v_{1},\cdots,v_{\mathbf{s}},v_{1}^{\ast},\cdots
,v_{s}^{\ast}\right\vert
\end{equation}%
\begin{align}
J_{A}\left\vert s_{1},\cdots,s_{s},s_{s+1},\cdots,s_{s+s}\right\rangle  &
=\left\vert s_{1},\cdots,s_{s},s_{s+1},\cdots,s_{s+s}\right\rangle \left(
\begin{array}
[c]{cc}%
\mathbb{O} & \mathbb{I}_{s}\\
-\mathbb{I}_{s} & \mathbb{O}%
\end{array}
\right)  =\left\vert -s_{s+1},\cdots,-s_{s+s},s_{1},\cdots,s_{s},\right\rangle
\\
\left\vert s_{1},\cdots,s_{s},s_{s+1},\cdots,s_{s+s}\right\rangle  &
=\left\vert s_{1},\cdots,s_{s},J_{A}s_{1},\cdots,J_{A}s_{s}\right\rangle
\end{align}%
\begin{equation}
J_{A}J_{A}^{t}=I_{s}=J_{A}^{t}J_{A}=-J_{A}^{2}=-I_{2s}%
\end{equation}%
\begin{equation}
\sigma_{A}=\left\vert s_{1},\cdots,s_{s},s_{s+1},\cdots,s_{s+s}\right\rangle
\left(
\begin{array}
[c]{cc}%
\mathbb{I}_{s} & \mathbb{O}\\
\mathbb{O} & -\mathbb{I}_{s}%
\end{array}
\right)  \left\langle s_{1},\cdots,s_{s},s_{s+1},\cdots,s_{s+s}\right\vert
\end{equation}%
\begin{equation}
\sigma_{A}\left\vert s_{1},\cdots,s_{s},s_{s+1},\cdots,s_{s+s}\right\rangle
\left(
\begin{array}
[c]{cc}%
\mathbb{I}_{s} & \mathbb{O}\\
\mathbb{O} & -\mathbb{I}_{s}%
\end{array}
\right)  =\left\vert s_{1},\cdots,s_{s},-s_{s+1},\cdots,-s_{s+s}\right\rangle
=\left\vert s_{1},\cdots,s_{s},-J_{A}s_{1},\cdots,-J_{A}s_{s}\right\rangle
\end{equation}%
\begin{equation}
\sigma_{A}=P_{A+}-P_{A-}%
\end{equation}%
\begin{equation}
\sigma_{A}^{2}=I_{2s}%
\end{equation}


\bigskip%
\begin{equation}
v^{\ast}=\sigma_{A}\left(  v\right)
\end{equation}%
\begin{equation}
V_{\mathbb{R}}=\left\{  v\left.  \sigma_{A}\left(  v\right)  \right\vert
=v,v\in V_{\mathbb{C}}\right\}
\end{equation}


\bigskip%

\begin{align}
\left\vert \mathbf{s}_{j}\right\rangle  &  =\frac{1}{\sqrt{2}}\left(
\left\vert \mathbf{v}_{j}\right\rangle +\left\vert \mathbf{v}_{j}^{\ast
}\right\rangle \right)  =\sqrt{2}\operatorname{Re}\left\vert \mathbf{v}%
_{j}\right\rangle \nonumber\\
\left\vert \mathbf{s}_{j+\mathbf{s}}\right\rangle  &  =\frac{-i}{\sqrt{2}%
}\left(  \left(  \left\vert \mathbf{v}_{j}\right\rangle -\left\vert
\mathbf{v}_{j}^{\ast}\right\rangle \right)  \right)  =\sqrt{2}%
\operatorname{Im}\left\vert \mathbf{v}_{j}\right\rangle ,
\end{align}%
\begin{equation}
\left[  J_{A},A\right]  =0
\end{equation}%
\begin{equation}
\left(
\begin{array}
[c]{cc}%
-\mathbb{I}_{s} & \mathbb{O}\\
\mathbb{O} & \mathbb{I}_{s}%
\end{array}
\right)  \left(
\begin{array}
[c]{cc}%
\mathbb{O} & \boldsymbol{\alpha}\\
-\boldsymbol{\alpha} & \mathbb{O}%
\end{array}
\right)  \left(
\begin{array}
[c]{cc}%
-\mathbb{I}_{s} & \mathbb{O}\\
\mathbb{O} & \mathbb{I}_{s}%
\end{array}
\right)  =\left(
\begin{array}
[c]{cc}%
-\mathbb{I}_{s} & \mathbb{O}\\
\mathbb{O} & \mathbb{I}_{s}%
\end{array}
\right)  \left(
\begin{array}
[c]{cc}%
\mathbb{O} & \boldsymbol{\alpha}\\
\boldsymbol{\alpha} & \mathbb{O}%
\end{array}
\right)  =\left(
\begin{array}
[c]{cc}%
\mathbb{O} & -\boldsymbol{\alpha}\\
\boldsymbol{\alpha} & \mathbb{O}%
\end{array}
\right)
\end{equation}%
\begin{equation}
\left(
\begin{array}
[c]{cc}%
-\mathbb{I}_{s} & \mathbb{O}\\
\mathbb{O} & \mathbb{I}_{s}%
\end{array}
\right)  \left(
\begin{array}
[c]{cc}%
-\boldsymbol{\alpha} & \mathbb{O}\\
\mathbb{O} & \boldsymbol{\alpha}%
\end{array}
\right)  \left(
\begin{array}
[c]{cc}%
-\mathbb{I}_{s} & \mathbb{O}\\
\mathbb{O} & \mathbb{I}_{s}%
\end{array}
\right)  =\left(
\begin{array}
[c]{cc}%
-\boldsymbol{\alpha} & \mathbb{O}\\
\mathbb{O} & \boldsymbol{\alpha}%
\end{array}
\right)
\end{equation}%
\begin{equation}
\sigma\left(  \operatorname{Re}x+i\operatorname{Im}x\right)
=\operatorname{Re}x-i\operatorname{Im}x
\end{equation}



\end{document}